% !TeX program = xelatex
\documentclass[a4paper,12pt]{article}

% ===== 1. KHỔ GIẤY, LỀ TRANG, FONT =====
\usepackage[left=3cm,right=2cm,top=2cm,bottom=2cm,footskip=1cm]{geometry}
\usepackage{fontspec}

\IfFontExistsTF{Times New Roman}{
  \setmainfont{Times New Roman}
}{
  \IfFileExists{fonts/times.ttf}{
    \setmainfont{times.ttf}[
      Path=fonts/,
      BoldFont=timesbd.ttf,
      ItalicFont=timesi.ttf,
      BoldItalicFont=timesbi.ttf]
  }{
    \IfFontExistsTF{TeX Gyre Termes}{
      \setmainfont{TeX Gyre Termes}
      \PackageWarningNoLine{MauBaoCao}{Times New Roman unavailable; using TeX Gyre Termes}
    }{
      \setmainfont{DejaVu Serif}
      \PackageWarningNoLine{MauBaoCao}{Times New Roman unavailable; using DejaVu Serif}
    }
  }
}
\usepackage{indentfirst}
\usepackage{microtype}
\usepackage{colortbl}
\usepackage{needspace}
\usepackage{graphicx}
\usepackage{float}
\usepackage{tikz}
\usepackage{amsmath,amssymb}
\usepackage{array,tabularx,booktabs,longtable}
\usepackage{enumitem}
\usepackage{titlesec}
\usepackage{tocloft}
\usepackage{caption}
\usepackage{xurl}
\usepackage[unicode,hidelinks]{hyperref}
\urlstyle{same}

% Nội dung cỡ 13 như trong Word (bp là điểm PDF), giãn dòng 1,5.
\renewcommand{\normalsize}{\fontsize{13bp}{19.5bp}\selectfont}
\AtBeginDocument{\normalsize}
\setlength{\parindent}{0.5cm}
\setlength{\parskip}{4pt}
\setlength{\emergencystretch}{3em}
\hyphenpenalty=10000
\exhyphenpenalty=10000
\raggedbottom
\setlist{topsep=4pt,itemsep=2pt,parsep=0pt,leftmargin=*}

% ===== 2. BA CẤP TIÊU ĐỀ: 1 / 1.1 / 1.1.1 =====
\setcounter{secnumdepth}{3}
\setcounter{tocdepth}{3}
\renewcommand{\thesection}{\arabic{section}}
\renewcommand{\thesubsection}{\thesection.\arabic{subsection}}
\renewcommand{\thesubsubsection}{\thesubsection.\arabic{subsubsection}}
\titleformat{\section}[hang]
  {\normalfont\fontsize{16}{20}\selectfont\bfseries}
  {\thesection.}{0.6em}{}
\titleformat{\subsection}[hang]
  {\normalfont\fontsize{14}{18}\selectfont\bfseries}
  {\thesubsection}{0.6em}{}
\titleformat{\subsubsection}[hang]
  {\normalfont\fontsize{13}{17}\selectfont\bfseries}
  {\thesubsubsection}{0.6em}{}
\titlespacing*{\section}{0pt}{16pt}{8pt}
\titlespacing*{\subsection}{0pt}{12pt}{6pt}
\titlespacing*{\subsubsection}{0pt}{8pt}{4pt}

% ===== 3. MỤC LỤC VÀ SỐ TRANG =====
\renewcommand{\contentsname}{MỤC LỤC}
\renewcommand{\listfigurename}{Danh sách hình vẽ}
\renewcommand{\listtablename}{Danh sách bảng}
\renewcommand{\cftloftitlefont}{\fontsize{16}{20}\selectfont\bfseries}
\renewcommand{\cftlottitlefont}{\fontsize{16}{20}\selectfont\bfseries}
\setlength{\cftbeforeloftitleskip}{0pt}\setlength{\cftafterloftitleskip}{14pt}
\setlength{\cftbeforelottitleskip}{0pt}\setlength{\cftafterlottitleskip}{14pt}
\renewcommand{\cftfigpresnum}{Hình }\renewcommand{\cfttabpresnum}{Bảng }
\renewcommand{\cftfigaftersnum}{.}\renewcommand{\cfttabaftersnum}{.}
\cftsetindents{figure}{0em}{5em}\cftsetindents{table}{0em}{5.2em}
\setlength{\cftbeforefigskip}{0pt}\setlength{\cftbeforetabskip}{0pt}
\renewcommand{\refname}{Tài liệu tham khảo}
\renewcommand{\cfttoctitlefont}{\fontsize{16}{20}\selectfont\bfseries}
\setlength{\cftbeforetoctitleskip}{0pt}
\setlength{\cftaftertoctitleskip}{18pt}
\renewcommand{\cftsecfont}{\normalfont\bfseries}
\renewcommand{\cftsecpagefont}{\normalfont\bfseries}
\renewcommand{\cftsubsecfont}{\normalfont}
\renewcommand{\cftsubsubsecfont}{\normalfont}
\renewcommand{\cftsecleader}{\cftdotfill{\cftdotsep}}
\cftsetindents{section}{0em}{2em}
\cftsetindents{subsection}{2em}{3em}
\cftsetindents{subsubsection}{5em}{4em}
\setlength{\cftbeforesecskip}{6pt}
\makeatletter
\def\ps@plain{%
  \let\@oddhead\@empty\let\@evenhead\@empty
  \def\@oddfoot{\hfil{\normalfont\fontsize{11}{13}\selectfont\thepage}\hfil}%
  \let\@evenfoot\@oddfoot}
\makeatother
\pagestyle{plain}

% ===== 4. BẢNG, HÌNH, CHÚ THÍCH =====
\counterwithin{figure}{section}
\counterwithin{table}{section}
\renewcommand{\topfraction}{0.9}\renewcommand{\bottomfraction}{0.8}
\renewcommand{\textfraction}{0.08}\renewcommand{\floatpagefraction}{0.75}
\setcounter{topnumber}{3}\setcounter{totalnumber}{4}
\renewcommand{\tablename}{Bảng}
\renewcommand{\figurename}{Hình}
\DeclareCaptionFont{reportcaption}{\fontsize{12}{15}\selectfont}
\captionsetup{font=reportcaption,labelfont=bf,labelsep=period,
  justification=centering,skip=6pt}
\newcommand{\bangfont}{\fontsize{11}{14}\selectfont\setlength{\tabcolsep}{4.2pt}\renewcommand{\arraystretch}{1.2}}
\newcommand{\bangnho}{\fontsize{10}{12.5}\selectfont\setlength{\tabcolsep}{4.2pt}\renewcommand{\arraystretch}{1.2}}
\newcolumntype{L}[1]{>{\raggedright\arraybackslash}m{#1}}
\newcolumntype{C}[1]{>{\centering\arraybackslash}m{#1}}
\newcommand{\dongdau}{\rowcolor[gray]{0.92}}
\newcommand{\code}[1]{\texttt{#1}}

% Hình giữ chỗ dùng \fbox, đoạn mã dùng listings (chỉ viết ASCII trong đoạn mã).
\usepackage{listings}
\lstset{
  basicstyle=\ttfamily\fontsize{9.5}{12.5}\selectfont,
  breaklines=true,
  columns=fullflexible,
  keepspaces=true,
  showstringspaces=false,
  frame=single,
  rulecolor=\color[gray]{0.6},
  captionpos=t,
  aboveskip=8pt,
  belowskip=8pt
}
\renewcommand{\lstlistingname}{Đoạn mã}

% ===== 5. CÔNG CỤ BIÊN SOẠN =====
% Bản nháp: \biensoantrue. Bản nộp: đổi thành \biensoanfalse để ẩn mọi ghi chú.
\newif\ifbiensoan
\biensoantrue
% \phutrach{TV2} : người viết mục
\newcommand{\phutrach}[1]{\ifbiensoan
  {\noindent\textbf{[Phụ trách: #1]}\par\vspace{2pt}}\fi}
% \ghichu{Nội dung ... Nguồn: ...}
\newcommand{\ghichu}[1]{\ifbiensoan
  {\itshape\noindent\textbf{Ghi chú biên soạn:} #1\par\vspace{6pt}}\fi}
% \hinhcho{mô tả}{chú thích}{nhãn} : khung giữ chỗ, thay bằng \hinh khi có ảnh
\newcommand{\hinhcho}[3]{%
  \begin{figure}[!htp]
  \centering
  \fbox{\parbox{0.92\textwidth}{\centering\small\itshape #1}}
  \caption{#2}\label{#3}
  \end{figure}}
% \hinh{tên file trong figures/}{bề rộng}{chú thích}{nhãn}
\newcommand{\hinh}[4]{%
  \begin{figure}[!htp]
  \centering
  \includegraphics[width=#2\textwidth]{figures/#1}
  \caption{#3}\label{#4}
  \end{figure}}
% Ô chưa có số liệu
\newcommand{\chuaco}{\dots}
% Tiêu đề Phần 1, Phần 2
\newcommand{\phanbaocao}[2]{%
  \clearpage
  \phantomsection
  \addcontentsline{toc}{section}{#1. #2}
  \begin{center}
    {\fontsize{16}{21}\selectfont\bfseries #1\par}
    \vspace{4pt}
    {\fontsize{14}{19}\selectfont\bfseries #2\par}
  \end{center}
  \vspace{6pt}}

% ===== 6. THÔNG TIN TRANG BÌA =====
\newcommand{\DonViChuQuan}{ĐẠI HỌC QUỐC GIA HÀ NỘI}
\newcommand{\TenTruong}{TRƯỜNG ĐẠI HỌC CÔNG NGHỆ}
\newcommand{\LoaiBaoCao}{BÁO CÁO BÀI TẬP LỚN}
\newcommand{\TenMonHoc}{HỆ QUẢN TRỊ CƠ SỞ DỮ LIỆU}
\newcommand{\TenDeTai}{CHI PHÍ CỦA CHỈ MỤC PHỤ TRÊN HỆ LƯU TRỮ KEY-VALUE PHÂN TÁN: NGHIÊN CỨU THỰC NGHIỆM TiDB SO VỚI MySQL InnoDB TRÊN BÀI TOÁN TRA CỨU VẬN ĐƠN}
\newcommand{\GiangVien}{[Họ tên giảng viên]}
\newcommand{\SoNhom}{Nhóm 7}
\newcommand{\LopMonHoc}{[Mã lớp môn học]}
\newcommand{\TenTVMot}{[Họ tên TV1]}
\newcommand{\MSSVMot}{[MSSV]}
\newcommand{\TenTVHai}{[Họ tên TV2]}
\newcommand{\MSSVHai}{[MSSV]}
\newcommand{\TenTVBa}{[Họ tên TV3]}
\newcommand{\MSSVBa}{[MSSV]}
\newcommand{\DiaDiemThoiGian}{Hà Nội, tháng 11 năm 2026}
\hypersetup{pdftitle={Chi phí của chỉ mục phụ: TiDB so với MySQL InnoDB},pdfauthor={Nhóm 7}}

\begin{document}

% ==============================================================================
% TRANG BÌA CHUẨN ĐỒ ÁN / BÀI TẬP LỚN ĐẠI HỌC CÔNG NGHỆ - ĐHQGHN
% ==============================================================================
\begin{titlepage}
\thispagestyle{empty}

% Khung viền đôi chuẩn đồ án UET
\begin{tikzpicture}[remember picture,overlay]
  \draw[line width=2pt]
    ([xshift=3cm,yshift=-2cm]current page.north west)
    rectangle ([xshift=-2cm,yshift=2cm]current page.south east);
  \draw[line width=0.5pt]
    ([xshift=3.12cm,yshift=-2.12cm]current page.north west)
    rectangle ([xshift=-2.12cm,yshift=2.12cm]current page.south east);
\end{tikzpicture}\par

\centering
\setlength{\parindent}{0pt}
\setlength{\parskip}{0pt}
\vspace*{0.5cm}

{\fontsize{13.5}{17}\selectfont\bfseries\DonViChuQuan\par}
\vspace{0.15cm}
{\fontsize{14.5}{19}\selectfont\bfseries\TenTruong\par}
\vspace{0.25cm}
\rule{5.5cm}{0.5pt}\par
\vspace{0.6cm}

% Logo trường ĐH Công nghệ
\IfFileExists{logo.png}{
  \includegraphics[width=3.6cm,height=3.6cm,keepaspectratio]{logo.png}
}{
  \fbox{\parbox[c][3.4cm][c]{3.4cm}{\centering\fontsize{11}{14}\selectfont LOGO TRƯỜNG}}
}\par
\vspace{0.6cm}

{\fontsize{19}{24}\selectfont\bfseries\LoaiBaoCao\par}
\vspace{0.3cm}
\begin{minipage}{0.9\textwidth}\centering
  {\fontsize{14}{18}\selectfont\bfseries MÔN HỌC: \TenMonHoc\par}
\end{minipage}\par
\vspace{0.5cm}

{\fontsize{12}{15}\selectfont\bfseries ĐỀ TÀI\par}
\vspace{0.15cm}
\begin{minipage}{0.92\textwidth}\centering
  {\fontsize{15}{20}\selectfont\bfseries\TenDeTai\par}
  \vspace{0.2cm}
  {\fontsize{11.5}{15}\selectfont\itshape\TenDeTaiTiengAnh\par}
\end{minipage}\par

\vspace{0.8cm}
{\fontsize{12.5}{18}\selectfont
\renewcommand{\arraystretch}{1.18}
\begin{tabular}{@{}>{\raggedright\arraybackslash\bfseries}p{4.8cm}>{\raggedright\arraybackslash}p{8.5cm}@{}}
  Giảng viên hướng dẫn: & \GiangVien\\
  Lớp môn học:          & \LopMonHoc\\
  Nhóm thực hiện:       & \SoNhom\\
  Sinh viên thực hiện:  & \textbf{Thành viên 1} -- MSSV: \dots\dots\dots\dots\ \textbf{(Trưởng nhóm)}\\
                        & \textit{(Phụ trách: Lý thuyết Kiến trúc \& Mô hình Toán học)}\\[2pt]
                        & \textbf{Thành viên 2} -- MSSV: \dots\dots\dots\dots\\
                        & \textit{(Phụ trách: Kỹ sư Hạ tầng, Cài đặt \& Thực nghiệm Index)}\\[2pt]
                        & \textbf{Thành viên 3} -- MSSV: \dots\dots\dots\dots\\
                        & \textit{(Phụ trách: Kỹ sư Benchmark, Chaos Test \& Đồ thị)}
\end{tabular}\par}

\vfill
{\fontsize{12.5}{16}\selectfont\bfseries\DiaDiemThoiGian\par}
\vspace*{0.3cm}
\end{titlepage}


% ===== 7. PHẦN MỞ ĐẦU (SỐ TRANG LA MÃ) =====
\pagenumbering{roman}
% Tóm tắt, Phân chia công việc. TV1 ghép, cả nhóm cùng sửa.
\section*{Tóm tắt}
\addcontentsline{toc}{section}{Tóm tắt}
\phutrach{TV1, viết sau cùng}
Một hệ thống quản lý vận đơn phải ghi liên tục, đồng thời tra cứu theo nhiều thuộc tính, nên bảng vận đơn thường mang nhiều chỉ mục phụ. MySQL InnoDB lưu mỗi chỉ mục phụ thành một cây B+Tree trên một máy; TiDB biến mỗi mục chỉ mục thành một cặp khóa--giá trị riêng, phân tán trên các node TiKV và nhân bản bằng Raft. Báo cáo này đo chi phí của chỉ mục phụ trên hai hệ ở ba mặt: dung lượng, độ trễ đọc và thông lượng ghi.

TiDB v8.5.5 (1 PD, 1 TiDB server, 3 TiKV) và MySQL 8.4.11 chạy luân phiên trong Docker trên cùng một máy, với cùng hồ sơ tài nguyên XS: mỗi hệ 4 CPU và khoảng 750 MiB bộ nhớ đệm dữ liệu. Dữ liệu vận đơn 1 triệu và 3 triệu dòng được sinh giống nhau trên ba máy từ cùng một giá trị seed. TiDB được đo với 3 bản sao mỗi Region (C3) và, để tách chi phí nhân bản, với 1 bản sao (C1). Mỗi phép đo lặp lại nhiều vòng; báo cáo dùng trung vị kèm khoảng dao động và không loại lượt nào vì kết quả trái giả thuyết.

Về dung lượng, mô hình dựng từ định dạng lưu trữ khớp số đo trong khoảng 1,2\% với MySQL và tới từng byte với TiDB, trừ chỉ mục UNIQUE; mỗi mục chỉ mục của TiDB lớn hơn MySQL 3,4 đến 4,8 lần, tính cho một bản sao, trước khi nén. Về đọc, ở 3 triệu dòng covering index giúp TiDB nhiều hơn khi ít dòng khớp (nhanh gấp 2,01 lần so với 1,08 lần trên MySQL khi $R = 10$), nhưng giúp MySQL nhiều hơn khi nhiều dòng khớp (5,91 so với 3,31 lần khi $R = 10.000$). Quét bảng bắt đầu nhanh hơn dùng chỉ mục ở khoảng 3,2\% số dòng trên TiDB và 6,9\% trên MySQL. Dùng mã vận đơn làm khóa chính làm chỉ mục phụ lớn thêm 27,4\% trên MySQL và 39,5\% trên TiDB. Về ghi, thêm 5 chỉ mục làm thông lượng thêm đơn giảm 18,9\% trên MySQL nhưng 82,7\% trên TiDB C3, và TiDB C1 không nhân bản vẫn giảm 79,4\%. Như vậy giả thuyết ``LSM-tree giúp TiDB giảm ít hơn'' bị bác bỏ, và nhân bản không phải nguyên nhân chính. Với covering index, cái giá khi ghi phụ thuộc vào việc cột thêm vào có bị cập nhật hay không: khi \code{status} bị cập nhật, thông lượng giảm 29,5\% trên MySQL và 61,6\% trên TiDB C3.

Các kết quả này có giới hạn rõ. Mọi node chạy trên một máy và một ổ đĩa, nên không có độ trễ mạng vật lý và chênh lệch giữa hai hệ không được quy cho mạng; khả năng chịu lỗi của C3 chỉ là chịu mất một tiến trình TiKV. Quy mô dừng ở 3 triệu dòng, con số 10 triệu dòng chỉ là ngoại suy từ mô hình; thí nghiệm ghi chỉ dùng một mức 64 luồng; và số tuyệt đối chỉ được so giữa hai hệ trên cùng một máy.

\clearpage
\section*{Phân chia công việc và mức độ đóng góp}
\renewcommand{\thetable}{P.\arabic{table}}
\addcontentsline{toc}{section}{Phân chia công việc và mức độ đóng góp}
\ghichu{Mỗi người sửa dòng của mình theo việc đã làm thật (đối chiếu lịch sử Git). Nhóm thống nhất tỉ lệ đóng góp trước khi nộp.}

\needspace{9\baselineskip}
{\bangfont
\begin{longtable}{|>{\raggedright\arraybackslash\bfseries}p{2.9cm}|>{\raggedright\arraybackslash}p{3.0cm}|>{\raggedright\arraybackslash}p{7.2cm}|>{\centering\arraybackslash}p{1.5cm}|}
\caption{Phân chia công việc và mức độ đóng góp của các thành viên}\label{tab:p1}\\
\hline
\dongdau \textbf{Thành viên} & \textbf{Vai trò} & \textbf{Công việc cụ thể} & \textbf{Tỉ lệ} \\
\hline
\endfirsthead
\multicolumn{4}{c}{\textbf{Bảng P.1.} Phân chia công việc và mức độ đóng góp (tiếp theo)}\\[2pt]
\hline
\dongdau \textbf{Thành viên} & \textbf{Vai trò} & \textbf{Công việc cụ thể} & \textbf{Tỉ lệ} \\
\hline
\endhead
\endfoot
\endlastfoot
\TenTVMot & Trưởng nhóm, máy M1 & Mô hình và số đo dung lượng (B.1); RQ1 ở 3 triệu dòng; kiến trúc TiDB; ghép báo cáo, slide & 33,3\% \\ \hline
\TenTVHai & Môi trường, máy M2 & Cấu hình TiDB và MySQL, công cụ đo dùng chung; RQ4; RQ3; code RQ1; cài đặt và toàn bộ demo A.3; B.3 & 33,3\% \\ \hline
\TenTVBa & Chi phí đọc, máy M3 & RQ1 ở 1 triệu dòng; RQ2; thực nghiệm sơ bộ B.0; B.2 & 33,3\% \\ \hline
\end{longtable}
}

\needspace{9\baselineskip}
{\bangfont
\begin{longtable}{|>{\raggedright\arraybackslash\bfseries}m{3.6cm}|>{\raggedright\arraybackslash}m{11.6cm}|}
\caption{Chia việc viết báo cáo}\label{tab:p2}\\
\hline
\dongdau \textbf{Thành viên} & \textbf{Phần phụ trách} \\
\hline
\endfirsthead
\multicolumn{2}{c}{\textbf{Bảng P.2.} Chia việc viết báo cáo (tiếp theo)}\\[2pt]
\hline
\dongdau \textbf{Thành viên} & \textbf{Phần phụ trách} \\
\hline
\endhead
\endfoot
\endlastfoot
\TenTVMot & Tóm tắt; Chương 1; Chương 3 (vấn đề, mục tiêu, phạm vi, cơ chế); Chương 4 (dữ liệu và dung lượng); Chương 7 (giả thuyết, khuyến nghị); Kết luận \\ \hline
\TenTVHai & Chương 2; Chương 3 (phương pháp, kịch bản); Chương 4 (RQ3); Chương 6 \\ \hline
\TenTVBa & Chương 4 (thực nghiệm sơ bộ); Chương 5; Chương 7 (hạn chế) \\ \hline
\end{longtable}
}


\clearpage
{\setlength{\parskip}{0pt}\setlength{\cftbeforesecskip}{4pt}\fontsize{13bp}{17.5bp}\selectfont\tableofcontents}

\clearpage
\phantomsection
\addcontentsline{toc}{section}{\listfigurename}
{\setlength{\parskip}{0pt}\fontsize{13bp}{17.5bp}\selectfont\listoffigures}

\clearpage
\phantomsection
\addcontentsline{toc}{section}{\listtablename}
{\setlength{\parskip}{0pt}\fontsize{13bp}{17.5bp}\selectfont\listoftables}

\clearpage
\section*{Danh mục các từ viết tắt}
\addcontentsline{toc}{section}{Danh mục các từ viết tắt}

{\bangfont
\begin{longtable}{|>{\centering\arraybackslash}m{1.1cm}|>{\raggedright\arraybackslash\bfseries}m{2.5cm}|>{\raggedright\arraybackslash}m{5.3cm}|>{\raggedright\arraybackslash}m{5.7cm}|}
\hline
\dongdau \textbf{STT} & \textbf{Từ viết tắt} & \textbf{Tiếng Anh đầy đủ} & \textbf{Nghĩa tiếng Việt} \\
\hline
\endfirsthead
\multicolumn{4}{c}{{\bfseries Danh mục các từ viết tắt (tiếp theo)}}\\
\hline
\dongdau \textbf{STT} & \textbf{Từ viết tắt} & \textbf{Tiếng Anh đầy đủ} & \textbf{Nghĩa tiếng Việt} \\
\hline
\endhead
\endfoot
\endlastfoot
1 & 2PC & Two-Phase Commit & Giao thức commit hai pha \\ \hline
2 & ACID & Atomicity, Consistency, Isolation, Durability & Bốn tính chất của giao dịch \\ \hline
3 & C1 & --- & Cấu hình 3 TiKV, 1 bản sao mỗi Region \\ \hline
4 & C3 & --- & Cấu hình 3 TiKV, 3 bản sao mỗi Region \\ \hline
5 & CBO & Cost-Based Optimizer & Bộ tối ưu truy vấn dựa trên chi phí \\ \hline
6 & CF & Column Family & Nhóm cột vật lý trong RocksDB \\ \hline
7 & COD & Cash On Delivery & Tiền thu hộ khi giao hàng \\ \hline
8 & gRPC & gRPC Remote Procedure Call & Giao thức gọi thủ tục từ xa giữa TiDB và TiKV \\ \hline
9 & IQR & Interquartile Range & Khoảng tứ phân vị \\ \hline
10 & KV & Key-Value & Cặp khóa -- giá trị \\ \hline
11 & LSM-tree & Log-Structured Merge-tree & Cây gộp có cấu trúc nhật ký \\ \hline
12 & MVCC & Multi-Version Concurrency Control & Điều khiển tương tranh đa phiên bản \\ \hline
13 & p95, p99 & 95th, 99th percentile & Bách phân vị 95, 99 của độ trễ \\ \hline
14 & PD & Placement Driver & Thành phần điều phối cụm TiDB \\ \hline
15 & PK & Primary Key & Khóa chính \\ \hline
16 & R & --- & Số dòng khớp điều kiện của một truy vấn \\ \hline
17 & RPC & Remote Procedure Call & Lời gọi thủ tục từ xa \\ \hline
18 & RQ & Research Question & Câu hỏi nghiên cứu \\ \hline
19 & SST & Sorted String Table & File dữ liệu đã sắp xếp của RocksDB \\ \hline
20 & TPS & Transactions Per Second & Số giao dịch mỗi giây \\ \hline
21 & TSO & Timestamp Oracle & Dịch vụ cấp dấu thời gian toàn cục \\ \hline
22 & W1, W2 & --- & Tải thêm đơn mới, tải cập nhật trạng thái \\ \hline
23 & WAL & Write-Ahead Log & Nhật ký ghi trước \\ \hline
24 & XS & --- & Hồ sơ tài nguyên dùng chung của nhóm \\ \hline
\end{longtable}
}


% ===== 8. NỘI DUNG CHÍNH =====
\clearpage
\pagenumbering{arabic}
\renewcommand{\thetable}{\thesection.\arabic{table}}
\setcounter{page}{1}

\section{Đặt vấn đề}
\label{sec:dat-van-de}
% Chương 1. Phụ trách: TV1. Chuyển từ bản Word của TV1 (26/11/2026), đã rút gọn.
\phutrach{TV1}

\subsection{Bối cảnh}
Thương mại điện tử và dịch vụ giao nhận tạo ra lượng vận đơn rất lớn, phải lưu và cập nhật liên tục. Tại Việt Nam, theo Bộ Công Thương, doanh thu thương mại điện tử năm 2025 ước đạt khoảng 31 tỷ USD, tăng 25,5\% so với năm trước~\cite{vn_ecommerce_2025}. Sản lượng bưu chính năm 2024 đạt trên 3,4 tỷ bưu gửi, trong đó dịch vụ phục vụ thương mại điện tử tăng 30\% so với năm 2023~\cite{vn_buu_chinh_2024}.
\ghichu{TV1 kiểm tra lại hai nguồn này. Bản Word ghi ``3,2 tỷ bưu gửi, 2,4 tỷ kiện thương mại điện tử'' theo EBI 2025, nhưng không tìm được nguồn khớp; các số GMV bốn sàn (Metric, YouNet ECI) và Pitney Bowes được bỏ cho gọn.}

Trong hệ thống quản lý vận đơn, cùng một bảng lớn vừa nhận ghi liên tục vừa phục vụ tra cứu theo nhiều thuộc tính. Mỗi vận đơn được tạo kèm số tiền thu hộ (COD), rồi được cập nhật trạng thái nhiều lần qua các chặng. Đồng thời, hệ thống cần tra một đơn theo mã vận đơn (\code{tracking\_code}), lấy các đơn mới nhất của một cửa hàng theo (\code{sender\_id}, \code{created\_at}) và đếm đơn theo kho đích và trạng thái (\code{dest\_hub\_id}, \code{status}). Mỗi kiểu tra cứu cần một chỉ mục phụ, nên chi phí của chỉ mục phụ về dung lượng, độ trễ đọc và thông lượng ghi trở thành câu hỏi trung tâm khi thiết kế.

\subsection{Vấn đề nghiên cứu}
Khi dữ liệu và tải ghi vượt khả năng của một máy chủ RDBMS, một hướng thường được cân nhắc là chuyển sang cơ sở dữ liệu SQL phân tán (NewSQL), vừa giữ mô hình quan hệ vừa mở rộng theo số node. Hai hệ đại diện lưu chỉ mục phụ theo cách khác nhau về bản chất. MySQL InnoDB lưu bảng và từng chỉ mục phụ thành các cây B+Tree trên một máy~\cite{mysql_innodb}. TiDB mã hóa mỗi dòng và mỗi mục chỉ mục thành một cặp Key-Value, chia theo Region trên các node TiKV, lưu bằng LSM-tree trong RocksDB, nhân bản bằng Multi-Raft và commit bằng giao dịch kiểu Percolator~\cite{pingcap_mapping,percolator}.

Câu hỏi của đề tài là: trên cùng bài toán vận đơn, dung lượng, độ trễ đọc và thông lượng ghi qua chỉ mục phụ thay đổi thế nào giữa hai hệ? Thực nghiệm sơ bộ bằng sysbench (đầu Chương~\ref{sec:dung-luong}) cho thấy giả thuyết ban đầu của nhóm không được xác nhận. Đó là động lực để nhóm xây dựng bộ thực nghiệm có kiểm soát trên bảng \code{shipments}, tách riêng từng nguồn chi phí thay vì kết luận chung hệ nào tốt hơn.

\subsection{Câu hỏi nghiên cứu và giả thuyết}
Bảng~\ref{tab:cau-hoi} liệt kê các câu hỏi nghiên cứu. Giả thuyết được đặt trước khi đo và có thể bị số liệu bác bỏ.
\begin{table}[!htp]
\centering\bangnho
\caption{Câu hỏi nghiên cứu, giả thuyết và lý do có thể sai}\label{tab:cau-hoi}
\begin{tabular}{|C{1.2cm}|L{3.6cm}|L{4.7cm}|L{4.7cm}|}
\hline
\dongdau \textbf{Mã} & \textbf{Câu hỏi} & \textbf{Giả thuyết} & \textbf{Lý do có thể sai} \\ \hline
B.1 & Bảng và chỉ mục chiếm bao nhiêu dung lượng? & Mô hình lý thuyết dự đoán được số đo sau khi hiệu chỉnh & Hệ số nén và phiên bản MVCC chưa dọn làm số đo lệch \\ \hline
RQ1 & Covering index giảm độ trễ đọc bao nhiêu? & R nhỏ: lợi ích covering ở TiDB lớn hơn vì chi phí cố định; R lớn: hai hệ có thể gần nhau & Point\_Get rẻ; TiDB gom handle theo lô; MySQL đọc ngẫu nhiên nếu vượt bộ nhớ đệm \\ \hline
RQ2 & Ở độ chọn lọc nào quét bảng nhanh hơn dùng chỉ mục? & Điểm giao nhau của TiDB thấp hơn MySQL & TiDB lấy dòng theo lô; TiKV ít CPU nên quét bảng cũng đắt \\ \hline
RQ3 & Kiểu khóa chính ảnh hưởng thế nào tới chỉ mục phụ? & Khóa chính dạng chuỗi làm chỉ mục phụ lớn hơn trên cả hai hệ; TiDB nonclustered phải tra hai bước khi tìm theo \code{id} & Nén và mã hóa khóa của TiKV có thể làm chênh lệch nhỏ đi \\ \hline
RQ4 & Thêm chỉ mục làm TPS ghi giảm bao nhiêu? & (a) TiDB giảm ít hơn nhờ LSM-tree; (b) TiDB giảm nhiều hơn vì mỗi chỉ mục thêm cặp KV phải prewrite và commit qua Raft & Hai giả thuyết đối lập, do thực nghiệm phân xử \\ \hline
\end{tabular}
\end{table}

\subsection{Mục tiêu và cách đo}
Đề tài không dừng ở việc so sánh độ trễ hay thông lượng tổng mà tìm nguyên nhân kỹ thuật (Bảng~\ref{tab:muc-tieu}). Với chiều đọc, nhóm ép đường truy cập bằng hint, đo các truy vấn đường chuẩn (\code{SELECT 1}, tra theo khóa chính) và khớp số đo với mô hình $t = a + bR$ để tách độ chọn lọc, toán tử truy cập và chi phí gọi KV. Cấu hình bản sao được kiểm tra trước mỗi lượt đo TiDB. Covering index được đánh giá trên cả ba mặt: đọc (RQ1), ghi (RQ4) và dung lượng (B.1).
\begin{table}[!htp]
\centering\bangfont
\caption{Mục tiêu và cách đo}\label{tab:muc-tieu}
\begin{tabular}{|L{3.4cm}|L{5.6cm}|L{5.6cm}|}
\hline
\dongdau \textbf{Nhóm mục tiêu} & \textbf{Mục tiêu} & \textbf{Cách đo} \\ \hline
Tách nguồn chênh lệch đọc & Tách độ chọn lọc, đường truy cập và chi phí gọi KV & Ép chỉ mục bằng hint, đường chuẩn, mô hình $t = a + bR$ \\ \hline
Cấu hình bản sao & Nêu đúng số TiKV, số bản sao, quorum và khả năng chịu lỗi & Kiểm tra trước khi đo TiDB \\ \hline
Covering trên truy vấn vận đơn & Kế hoạch, độ trễ, chi phí ghi, dung lượng & RQ1, RQ4 và B.1 \\ \hline
\end{tabular}
\end{table}

\subsection{Phạm vi và giới hạn}
Đề tài so sánh TiDB v8.5.5 với MySQL 8.4.11, chạy trong Docker trên WSL2 theo hồ sơ tài nguyên chung XS: mỗi hệ dùng tổng cộng 4 CPU; bộ nhớ đệm dữ liệu là 750 MiB ở TiDB (3 TiKV, mỗi node 250 MiB block cache) và 752 MiB ở MySQL (buffer pool). Dữ liệu có hai mức 1 triệu và 3 triệu dòng; mức 10 triệu dòng chỉ được ngoại suy bằng mô hình dung lượng ở Chương~\ref{sec:dung-luong}, không đo. RQ4 chỉ đo ở 1 triệu dòng với 64 luồng; lý do bỏ mức 16 luồng nêu ở Chương~\ref{sec:ghi}.

RQ3 chỉ đo ảnh hưởng của kiểu khóa chính tới dung lượng chỉ mục phụ và tới việc tra cứu, ở 1 triệu dòng trên một máy. Các thí nghiệm còn lại đều giữ cùng lược đồ \code{PRIMARY KEY (id)} trên hai hệ để cô lập chi phí của riêng chỉ mục phụ; vì vậy các biến thể giảm điểm nóng ghi như \code{AUTO\_RANDOM} của TiDB không được thử.

So với đề cương đã đăng ký, nhóm có hai điều chỉnh. Thứ nhất, đề cương dự kiến 1 đến 10 triệu bản ghi; với hồ sơ XS trên laptop 16 GB, nhóm chỉ đo ở 1 và 3 triệu dòng, còn 10 triệu dòng chỉ được tính bằng mô hình dung lượng. Thứ hai, đề cương dự kiến so sánh độ trễ mạng phân tán ở RQ1; theo nhận xét của giảng viên, nhóm tách riêng chi phí gọi KV và không quy chênh lệch cho mạng, vì mọi node nằm trên một máy.

Mọi container của cụm TiDB (1 PD, 1 TiDB server, 3 TiKV) và MySQL chạy luân phiên trên cùng một máy và một ổ đĩa, nên giữa các node không có độ trễ mạng vật lý và chênh lệch hiệu năng không thể quy cho đường truyền. Khối lượng đo chia trên ba máy (Bảng~\ref{tab:ba-may}); số tuyệt đối chỉ được so giữa hai hệ trên cùng một máy, còn giữa các máy chỉ đối chiếu tỉ lệ. Demo tắt node và thêm node ở Chương~\ref{sec:cai-dat} chỉ minh họa cách cụm vận hành, không dùng để kết luận về chi phí chỉ mục.

\subsection{Cấu trúc báo cáo}
Phần 1 gồm Chương~\ref{sec:kien-truc} (tổng quan và kiến trúc TiDB) và Chương~\ref{sec:cai-dat} (cài đặt, cấu hình, vận hành và demo). Phần 2 gồm Chương~\ref{sec:phuong-phap} (vấn đề, cơ chế và phương pháp), Chương~\ref{sec:dung-luong} (dữ liệu, dung lượng và RQ3), Chương~\ref{sec:doc} (chi phí đọc), Chương~\ref{sec:ghi} (chi phí ghi) và Chương~\ref{sec:tong-hop} (tổng hợp). Lệnh, cấu hình và mã nguồn nằm trong kho mã của nhóm, dẫn ở đầu Chương~\ref{sec:cai-dat}.


\phanbaocao{PHẦN 1}{Giới thiệu và cài đặt hệ quản trị cơ sở dữ liệu TiDB}

\section{Tổng quan và kiến trúc TiDB}
\label{sec:kien-truc}
% Mục A.1. Phụ trách: TV1. Chuyển từ bản Word của TV1 (26/11/2026), đã rút gọn.
\phutrach{TV1 -- mục A.1}

\subsection{Đặc trưng dữ liệu vận đơn}
Bảng \code{shipments} của đề tài mô phỏng một hệ thống vận đơn với 100.000 cửa hàng (số đơn mỗi cửa hàng theo phân phối Zipf, tham số 1,1), 200 kho trung chuyển và 80\% đơn ở trạng thái đã giao. Trên cùng bảng này, hệ thống phải đáp ứng đồng thời ba nhóm yêu cầu:
\begin{itemize}
  \item \textbf{Giao dịch ACID} khi tạo đơn, cập nhật trạng thái và điều chỉnh số tiền thu hộ (\code{cod\_amount}).
  \item \textbf{Tra cứu qua nhiều chỉ mục phụ}: tra điểm theo \code{tracking\_code}, tra theo cửa hàng và thời gian tạo (\code{sender\_id}, \code{created\_at}), lọc theo khoảng \code{cod\_amount}.
  \item \textbf{Thống kê theo kho và trạng thái} (\code{dest\_hub\_id}, \code{status}), chạy song song với luồng cập nhật.
\end{itemize}

\subsection{Các nhóm hệ quản trị và hạn chế}
\textbf{RDBMS trên một máy.} MySQL (InnoDB) hay PostgreSQL hỗ trợ đầy đủ giao dịch ACID, SQL chuẩn và nhiều loại chỉ mục phụ, kể cả covering index trên B+Tree. Tuy nhiên toàn bộ bảng và chỉ mục nằm trên một node, nên dung lượng, bộ nhớ đệm và khả năng chịu lỗi bị giới hạn bởi phần cứng của máy đó. Khi tải ghi vượt một máy, phải chia dữ liệu thủ công (sharding) ở tầng ứng dụng hoặc tầng trung gian, khiến giao dịch và chỉ mục toàn cục phức tạp hơn.

\textbf{Wide-column NoSQL.} Cassandra hay HBase ghi nhanh và mở rộng ngang tốt. Đổi lại, nhóm này không có giao dịch ACID nhiều dòng, nhiều bảng theo mô hình quan hệ, chỉ mục phụ hạn chế và không hỗ trợ JOIN, nên khó phục vụ các truy vấn lọc theo nhiều thuộc tính độc lập của bài toán vận đơn.

\textbf{NewSQL phân tán.} TiDB giữ mô hình quan hệ, giao thức MySQL và giao dịch ACID ở tầng trên, còn tầng lưu trữ là các cặp Key-Value tự động chia thành Region trên nhiều node TiKV~\cite{pingcap_arch}. Bảng~\ref{tab:nhom-hqt} đối chiếu ba nhóm với yêu cầu của bài toán vận đơn. Cái giá của cách tổ chức này đối với chỉ mục phụ là câu hỏi mà đề tài đo.

\begin{table}[!htp]
\centering\bangfont
\caption{So sánh các nhóm hệ quản trị với bài toán vận đơn}\label{tab:nhom-hqt}
\begin{tabular}{|L{2.6cm}|L{3.9cm}|L{3.9cm}|L{3.9cm}|}
\hline
\dongdau \textbf{Tiêu chí} & \textbf{RDBMS một máy} & \textbf{Wide-column} & \textbf{NewSQL (TiDB)} \\ \hline
Giao dịch ACID & Đầy đủ, trên một máy & Hạn chế giao dịch nhiều dòng, nhiều bảng & ACID phân tán nhiều dòng (Snapshot Isolation) \\ \hline
Chỉ mục phụ & Đầy đủ trên B+Tree, có covering index & Hạn chế, khó truy vấn theo nhiều thuộc tính & Có, kể cả covering index, trên nền KV phân tán \\ \hline
Mở rộng ngang & Giới hạn bởi một máy; ghi nhiều phải sharding & Tốt cho cả lưu trữ và tải ghi & Mở rộng trực tuyến, tách tầng SQL và tầng lưu trữ \\ \hline
Chịu lỗi & Phụ thuộc một máy nếu không dựng thêm bản sao & Nhờ bản sao trên nhiều node & Tự bầu leader khi mất thiểu số bản sao (Multi-Raft) \\ \hline
\end{tabular}
\end{table}

% Tổng quan TiDB. Phụ trách: TV1.
\phutrach{TV1 -- tổng quan TiDB}

\subsection{Tổng quan về TiDB}
\ghichu{Khoảng nửa trang, trình bày tổng quát: TiDB là gì, ai phát triển, mã nguồn mở theo giấy phép nào, phiên bản v8.5.5, tương thích giao thức MySQL, nhóm bài toán phù hợp và không phù hợp. Mỗi thông tin có trích dẫn, ví dụ \cite{pingcap_arch}. Chi tiết dẫn sang Phụ lục C.}

\subsection{Các thành phần chính}
\begin{table}[!htp]
\centering\bangfont
\caption{Các thành phần chính của TiDB}\label{tab:thanh-phan}
\begin{tabular}{|L{2.6cm}|L{6.6cm}|C{2.4cm}|L{2.8cm}|}
\hline
\dongdau \textbf{Thành phần} & \textbf{Vai trò} & \textbf{Lưu trạng thái} & \textbf{Trong đề tài} \\ \hline
TiDB server & \chuaco & \chuaco & 1 node \\ \hline
TiKV & \chuaco & \chuaco & 3 node \\ \hline
PD & \chuaco & \chuaco & 1 node \\ \hline
TiFlash & \chuaco & \chuaco & Không dùng \\ \hline
TiDB Dashboard & \chuaco & --- & Có, trong PD \\ \hline
\end{tabular}
\end{table}

% Mục A.2. Phụ trách: TV1.
\phutrach{TV1 -- mục A.2; mỗi mục kết thúc bằng ảnh hưởng tới chỉ mục}

\subsection{Tách tầng tính toán và tầng lưu trữ}
\ghichu{TiDB server, TiKV, Coprocessor. Giải thích \code{TableReader}, \code{IndexReader} (covering), \code{IndexLookUp}, \code{Point\_Get}, vì Chương~\ref{sec:doc} đo chúng trực tiếp.}

\hinhcho{Sơ đồ 1 PD, 1 TiDB server, 3 TiKV đúng với cụm đã cài}{Kiến trúc cụm TiDB dùng trong đề tài}{fig:kien-truc}

\subsection{Ánh xạ bảng và chỉ mục sang Key-Value}

\subsubsection{Dòng dữ liệu và chỉ mục}
\ghichu{Dạng key của dòng, chỉ mục unique và không unique; đối chiếu với B+Tree phụ của InnoDB.}

\subsubsection{RocksDB và LSM-tree}
\ghichu{WAL, MemTable, SST, compaction. Nêu cả lợi ích (không sửa trang tại chỗ) và cái giá (khuếch đại ghi). Không kết luận bên nào rẻ hơn; RQ4 phân xử.}

\subsection{Region và Multi-Raft}
\ghichu{Kích thước Region theo \code{SHOW CONFIG} trên v8.5.5. Quorum là đa số peer của một Region. Định nghĩa C3 và C1 dùng Bảng~\ref{tab:ban-sao}. Ba TiKV là ba container trên một máy. Bỏ số ``3--8 giây'' ở bản cũ.}

\subsection{Placement Driver, TSO và giao dịch Percolator}
\ghichu{TSO, prewrite, commit, MVCC. Mức cô lập Snapshot Isolation, không hỗ trợ SERIALIZABLE. Số key phải ghi trong một giao dịch tăng theo số chỉ mục.}

\subsection{Kết luận chương}
\begin{table}[!htp]
\centering\bangfont
\caption{Cơ chế của TiDB và câu hỏi nghiên cứu liên quan}\label{tab:co-che}
\begin{tabular}{|L{4.6cm}|L{6.6cm}|C{3.0cm}|}
\hline
\dongdau \textbf{Cơ chế} & \textbf{Ảnh hưởng tới chỉ mục} & \textbf{Câu hỏi} \\ \hline
IndexReader, IndexLookUp & \chuaco & RQ1, RQ2 \\ \hline
Ánh xạ KV, LSM-tree & \chuaco & B.1, RQ4 \\ \hline
Region, Multi-Raft & \chuaco & RQ4 (C1, C3) \\ \hline
TSO, Percolator & \chuaco & RQ1, RQ4 \\ \hline
\end{tabular}
\end{table}


\clearpage
\section{Cài đặt, cấu hình, vận hành và demo}
\label{sec:cai-dat}
% Mục A.3 phần cài đặt. Phụ trách: TV2.
\phutrach{TV2 -- mục A.3, phần cài đặt}

Chương này chỉ trình bày môi trường và cách cài đặt ở mức tổng quát. Toàn bộ lệnh, file cấu hình, mã nguồn và kết quả đo nằm trong kho mã của nhóm (\url{https://github.com/taesikwoo268/nhom7-tidb-indexing}); muốn dựng lại môi trường, chỉ cần làm theo file \path{README.md}.

\subsection{Ba máy thực nghiệm}

Mỗi thành viên đo phần của mình trên máy cá nhân (Bảng~\ref{tab:ba-may}). Ba máy có phần cứng khác nhau, nên nhóm chỉ so số tuyệt đối giữa hai hệ trên cùng một máy; giữa các máy, nhóm chỉ so tỉ lệ.

\begin{table}[!htp]
\centering\bangnho
\caption{Ba máy thực nghiệm}\label{tab:ba-may}
\begin{tabular}{|C{1.1cm}|L{2.6cm}|L{3.2cm}|C{1.3cm}|L{2.4cm}|L{3.2cm}|}
\hline
\dongdau \textbf{Máy} & \textbf{Phần đo} & \textbf{CPU, số luồng} & \textbf{RAM} & \textbf{Ổ dữ liệu} & \textbf{Windows, Docker, WSL2} \\ \hline
M1 & RQ1 (3 triệu dòng, đối chiếu) & Intel Core i5-1235U, 10 nhân, 12 luồng & 16 GB & SSD NVMe 256 GB (KIOXIA) & Windows 11 Pro (10.0.26200); Docker 29.6.2; WSL2 10 GB RAM, 4 GB swap \\ \hline
M2 & RQ4, RQ3 (1 triệu dòng); RQ1, B.1 (1 và 3 triệu dòng) & Intel Core Ultra 5 125H, 14 nhân, 18 luồng & 16 GB & SSD NVMe 512 GB (UMIS) & Windows 11 Home (10.0.26300); Docker 29.8.0; WSL2 10 GB RAM, 4 GB swap \\ \hline
M3 & RQ1 (1 triệu dòng); RQ2 & Intel Core i5-13420H, 8 nhân, 12 luồng & 16 GB & SSD NVMe 512 GB (SK hynix) & Windows 11 Home Single Language (10.0.26200); Docker 29.8.2; WSL2 10 GB RAM, 4 GB swap \\ \hline
\end{tabular}
\end{table}

Với giới hạn WSL2 như trên, Docker thấy 12 CPU và 9,7 GiB bộ nhớ trên M1, 18 CPU và 9,7 GiB trên M2, 10 CPU và 9,7 GiB trên M3. Trong các lượt đo chính thức, máy luôn được cắm sạc để chế độ tiết kiệm pin không làm thay đổi hiệu năng.

\subsection{Cấu hình tài nguyên}

Cả ba máy dùng chung một hồ sơ tài nguyên mà nhóm gọi là XS (Bảng~\ref{tab:xs}). Hồ sơ này được thu nhỏ để cụm TiDB chạy được trên laptop 16 GB, trong khi hai hệ vẫn có cùng tổng giới hạn CPU và bộ nhớ đệm dữ liệu gần bằng nhau. Mỗi lúc chỉ có một hệ chạy.

\begin{table}[!htp]
\centering\bangfont
\caption{Hồ sơ XS: giới hạn tài nguyên của từng thành phần}\label{tab:xs}
\begin{tabular}{|L{4.2cm}|C{2.4cm}|C{3.0cm}|L{4.6cm}|}
\hline
\dongdau \textbf{Thành phần} & \textbf{CPU} & \textbf{RAM (MiB)} & \textbf{Bộ nhớ đệm} \\ \hline
PD & 0,4 & 512 & --- \\ \hline
TiDB server v8.5.5 & 1,2 & 1.280 & --- \\ \hline
TiKV v8.5.5 (3 node) & $3 \times 0{,}8$ & $3 \times 1.280$ & Block cache $3 \times 250$ MiB \\ \hline
\textbf{Tổng cụm TiDB} & \textbf{4,0} & \textbf{5.632} & \textbf{750 MiB} \\ \hline
MySQL 8.4.11 & 4,0 & 2.048 & Buffer pool 752 MiB (đặt 750M) \\ \hline
sysbench (RQ4) & 2,0 & 512 & --- \\ \hline
\end{tabular}
\end{table}

Hai hệ có cùng tổng 4 CPU và bộ nhớ đệm dữ liệu gần bằng nhau, nhưng tổng RAM khác nhau: cụm TiDB được 5,5 GiB, MySQL 2 GiB, vì cụm TiDB có năm tiến trình và tiến trình nào cũng cần bộ nhớ riêng. MySQL là một tiến trình nên dùng được cả 4 CPU, còn mỗi TiKV bị giới hạn ở 0,8 CPU kể cả khi hai TiKV kia rảnh; Chương~\ref{sec:ghi} cho thấy giới hạn này có ảnh hưởng thật. Buffer pool đặt 750M được MySQL làm tròn thành 752 MiB.

MySQL đặt \code{innodb\_flush\_log\_at\_trx\_commit=1} và tắt binlog; TiDB giữ cấu hình Raft log mặc định, mỗi TiKV có \code{memory-usage-limit} 1.000 MiB. Image TiKV chính thức được loại phần thông tin debug, mã thực thi giữ nguyên; nhờ vậy bộ nhớ một TiKV chiếm khi cụm rỗng giảm từ khoảng 1,9 GB xuống 0,16 GB. Về collation, MySQL dùng \code{utf8mb4\_unicode\_ci} còn TiDB dùng \code{utf8mb4\_bin}; chỉ \code{tracking\_code} là cột chuỗi. Mục~\ref{sec:rq4-collation} cho thấy khác biệt này không vượt mức nhiễu ở RQ4; RQ1 vẫn dùng collation mặc định nên ở Q1 chỉ so tỉ số trong từng hệ.

\subsection{Số node, số bản sao, quorum và khả năng chịu lỗi}

Số node TiKV và số bản sao của mỗi Region là hai khái niệm khác nhau, nên cần nói rõ cả hai. Trong đề tài, cả hai cấu hình TiDB đều có đúng 3 TiKV; nhóm chỉ đổi số bản sao qua tham số \code{replication.max-replicas} của PD (Bảng~\ref{tab:ban-sao}).

\begin{table}[!htp]
\centering\bangnho
\caption{Định nghĩa cấu hình bản sao}\label{tab:ban-sao}
\begin{tabular}{|C{1.4cm}|C{1.4cm}|C{2.0cm}|C{1.8cm}|C{1.7cm}|L{5.5cm}|}
\hline
\dongdau \textbf{Ký hiệu} & \textbf{TiKV} & \textbf{Bản sao mỗi Region} & \textbf{Quorum} & \textbf{Chịu mất} & \textbf{Dùng cho} \\ \hline
C3 & 3 & 3 & 2 trên 3 & 1 store & Mọi phép đo chính \\ \hline
C1 & 3 & 1 & 1 trên 1 & Không & Tách chi phí nhân bản ở RQ4, bộ 0 và 5 chỉ mục \\ \hline
\end{tabular}
\end{table}

Ở C3, một lệnh ghi được commit khi 2 trong 3 peer đã ghi Raft log~\cite{raft}, nên mất một store thì các Region vẫn đủ đa số để bầu leader mới; ở C1, store chứa peer duy nhất dừng thì Region cũng dừng. MySQL chạy một node, không có bản sao. Một TiKV duy nhất không thể tạo cấu hình 3 bản sao vì các peer của một Region phải nằm trên các store khác nhau. Ba TiKV ở đây là ba container trên một máy và một ổ đĩa, nên ``chịu mất một store'' chỉ có nghĩa là chịu mất một tiến trình; PD và TiDB server chỉ có một node nên vẫn là điểm lỗi đơn.

Trước mỗi lượt đo TiDB, chương trình kiểm tra \code{max-replicas}, đủ 3 store ở trạng thái \code{Up} và mọi Region có đúng số peer; thiếu một điều kiện là chương trình dừng.

\subsection{Các bước cài đặt}

Việc cài đặt gồm năm bước. Vì mọi cấu hình đều nằm trong Git, làm đúng các bước này thì ba máy sẽ có môi trường giống nhau.

\begin{enumerate}
  \item Cài Docker Desktop với WSL2, đặt \code{memory=10GB} và \code{swap=4GB} cho WSL2.
  \item Lấy mã nguồn về và tạo môi trường Python 3.10 trở lên.
  \item Khai báo tên máy và hồ sơ (\code{MACHINE\_ID}, \code{PROFILE=XS}). Nếu cổng 3306 đã bị chiếm thì đổi cổng MySQL; M2 dùng cổng 3307.
  \item Chạy lệnh kiểm tra. Lệnh này từ chối khi sai hồ sơ, Docker thiếu bộ nhớ hoặc CPU, máy đang chạy pin, hoặc có file cấu hình khác với Git.
  \item Bật từng hệ bằng \code{use\_engine.ps1} và nạp dữ liệu bằng \code{load.ps1}.
\end{enumerate}

\subsection{Minh chứng cài đặt}

Hình~\ref{fig:minh-chung-cai-dat} là ảnh chụp terminal trên M2 khi bật cụm TiDB: lệnh kiểm tra cho \code{team\_hash} chung của nhóm và kết quả \code{DAT}, năm container của cụm đang chạy, phiên bản v8.5.5 và ba store ở trạng thái Up với \code{max-replicas} bằng 3. Log nạp dữ liệu cho thấy hai hệ cùng nạp 1 triệu dòng với cùng tổng \code{cod\_amount} (5000968738439) và cùng mã SHA-256 của file nguồn; MySQL báo phiên bản 8.4.11, còn TiDB báo chuỗi phiên bản tương thích MySQL 8.0.11 kèm phiên bản thật v8.5.5.

\hinh{minh_chung_cai_dat.png}{0.8}{Ảnh chụp terminal trên M2: kiểm tra cấu hình, các container đang chạy, phiên bản TiDB và trạng thái ba store}{fig:minh-chung-cai-dat}

% Mục A.3 phần vận hành. Phụ trách: TV2.
\phutrach{TV2 -- mục A.3, phần vận hành}

\subsection{Vận hành cụm}
\ghichu{Trình bày tổng quát, mỗi việc một đến hai câu kèm lệnh; chi tiết dẫn sang Phụ lục C.}

\begin{table}[!htp]
\centering\bangnho
\caption{Các thao tác vận hành thường dùng}\label{tab:van-hanh}
\begin{tabular}{|L{4.4cm}|L{7.0cm}|L{3.2cm}|}
\hline
\dongdau \textbf{Thao tác} & \textbf{Lệnh} & \textbf{Kết quả mong đợi} \\ \hline
Bật TiDB, tắt MySQL & \code{.\textbackslash scripts\textbackslash use\_engine.ps1 -Engine tidb -Switch} & \chuaco \\ \hline
Bật MySQL, tắt TiDB & \code{.\textbackslash scripts\textbackslash use\_engine.ps1 -Engine mysql -Switch} & \chuaco \\ \hline
Xem dịch vụ đang chạy & \code{.\textbackslash scripts\textbackslash use\_engine.ps1 -Status} & \chuaco \\ \hline
Dừng tất cả, giữ dữ liệu & \code{.\textbackslash scripts\textbackslash use\_engine.ps1 -Stop} & \chuaco \\ \hline
Kiểm tra cấu hình chung & \code{python -m bench.cli check} & \chuaco \\ \hline
Nạp dữ liệu & \code{.\textbackslash scripts\textbackslash load.ps1 -Engine tidb -Rows 1000000} & \chuaco \\ \hline
\end{tabular}
\end{table}

\subsubsection{Kiểm tra trạng thái}
\ghichu{Trạng thái store và số bản sao của Region: \code{information\_schema.tikv\_store\_status}, \code{information\_schema.tikv\_region\_peers}, \code{SHOW CONFIG}. Trích kết quả thật trên M2.}

\subsubsection{Giám sát bằng TiDB Dashboard}
\ghichu{Dashboard ở cổng 2379 của PD: Cluster Info, Key Visualizer, SQL Statements. Ảnh chụp trên hồ sơ XS, hoặc ghi rõ hồ sơ và ngày chụp.}

\hinhcho{Ảnh TiDB Dashboard: 1 PD, 1 TiDB, 3 TiKV ở trạng thái Up}{TiDB Dashboard trên máy M2}{fig:dashboard}

\hinhcho{Ảnh Key Visualizer}{Key Visualizer trên máy M2}{fig:keyviz}

% Mục A.3 demo chức năng cơ bản. Phụ trách: TV2.
\phutrach{TV2 -- mục A.3, demo chức năng cơ bản}

\subsection{Demo chức năng cơ bản}
Demo chạy trên database riêng \code{demo\_nhom7} với 20.000 vận đơn, không dùng bảng của các thí nghiệm. Cùng một kịch bản chạy trên cả hai hệ.
\ghichu{Chạy \path{scripts/demo_co_ban.py} trên M2, đầu ra ở \path{results/runs}. Trích phần chính vào báo cáo, toàn bộ đầu ra dẫn sang Phụ lục C. Đây là demo chức năng, không phải phép đo hiệu năng.}

\begin{table}[!htp]
\centering\bangfont
\caption{Các chức năng được demo}\label{tab:demo-co-ban}
\begin{tabular}{|L{4.2cm}|L{5.6cm}|C{2.2cm}|C{2.2cm}|}
\hline
\dongdau \textbf{Chức năng} & \textbf{Kiểm tra} & \textbf{MySQL} & \textbf{TiDB} \\ \hline
Kết nối, phiên bản & Đọc được phiên bản, 3 store Up (TiDB) & \chuaco & \chuaco \\ \hline
Tạo bảng & \code{SHOW CREATE TABLE} & \chuaco & \chuaco \\ \hline
Thêm, đọc, sửa, xóa & Dòng thay đổi đúng như lệnh & \chuaco & \chuaco \\ \hline
Giao dịch ROLLBACK & Dữ liệu giữ giá trị cũ & \chuaco & \chuaco \\ \hline
Giao dịch COMMIT & Dữ liệu đổi sang giá trị mới & \chuaco & \chuaco \\ \hline
Ràng buộc UNIQUE & Từ chối mã vận đơn trùng & \chuaco & \chuaco \\ \hline
Tạo chỉ mục & Kế hoạch đổi sau khi tạo chỉ mục & \chuaco & \chuaco \\ \hline
\end{tabular}
\end{table}

\subsubsection{Thêm, đọc, sửa, xóa và giao dịch}
\ghichu{Trích lệnh và kết quả.}

\subsubsection{Chỉ mục và kế hoạch thực thi}
\ghichu{Trích \code{EXPLAIN} ba lần: chưa có chỉ mục, có chỉ mục không covering, có chỉ mục covering. Chỉ nêu toán tử, không nêu thời gian.}

% Mục A.3 phần demo bổ sung. Phụ trách: TV2 (nhận từ TV3 ngày 25/11/2026).
% Nguồn: results/demo_chiu_loi_{tidb,mysql}_M2.csv, results/demo_them_node_M2.csv, results/runs/demo_*_M2.txt.
\phutrach{TV2 -- mục A.3, phần demo bổ sung}

\subsection{Demo bổ sung: chịu lỗi}

Demo minh họa cơ chế quorum ở Bảng~\ref{tab:ban-sao}, không dùng cho kết luận nào về chỉ mục. Trên cụm C3 (3 TiKV, 3 bản sao), 8 luồng liên tục chèn từng dòng vào một bảng thử và ghi nhớ mọi giao dịch đã được xác nhận. Ở giây 30, container TiKV đang giữ leader của bảng thử bị tắt đột ngột bằng \code{docker compose kill}; ở giây 90 container được bật lại. Store bị tắt được chọn có chủ đích để sự cố chắc chắn ảnh hưởng tới tải. Trên MySQL, cùng kịch bản thời gian được áp dụng cho container duy nhất.

\begin{table}[!htp]
\centering\bangfont
\caption{Demo chịu lỗi trên máy M2, ngày 25/11/2026}\label{tab:demo-loi}
\begin{tabular}{|L{6.4cm}|C{4.0cm}|C{4.0cm}|}
\hline
\dongdau \textbf{Chỉ số} & \textbf{TiDB C3, tắt \code{tikv2}} & \textbf{MySQL một node} \\ \hline
Số giây không có giao dịch thành công & 13 & 61 \\ \hline
Thời gian về 90\% mức nền & 14 giây & 62 giây \\ \hline
Phục hồi nhờ & Bầu lại leader, tự động & Bật lại container ở giây 90 \\ \hline
Lệnh ghi bị lỗi & 0 / 129.275 & 296 / 423.772 \\ \hline
Giao dịch đã xác nhận bị mất & 0 / 129.275 & 0 / 423.476 \\ \hline
Dòng có trong bảng nhưng client không nhận được xác nhận & 0 & 7 \\ \hline
\end{tabular}
\end{table}

\hinh{demo_chiu_loi_tidb_M2.png}{0.9}{Giao dịch thành công mỗi giây khi tắt một TiKV ở cấu hình C3}{fig:demo-loi-tidb}

Ngay sau khi \code{tikv2} bị tắt, tải dừng hẳn vì Region của bảng thử mất leader (Hình~\ref{fig:demo-loi-tidb}). Hai peer còn lại vẫn đủ quorum 2 trên 3 nên bầu được leader mới; tải chạy lại ở giây 44 mà không cần can thiệp. Theo PD, số leader trên \code{tikv2} giảm từ 27 về 0, trong khi \code{tikv1} tăng từ 28 lên 40 và \code{tikv3} từ 19 lên 34. Các luồng ghi trên TiDB không gặp lỗi: câu lệnh bị treo trong lúc chờ leader mới rồi hoàn thành. Sau khi được bật lại, \code{tikv2} nhận lại leader, đạt 18 leader ở cuối lượt đo.

Với MySQL, 61 giây ngừng phục vụ chủ yếu là khoảng chờ 60 giây mà kịch bản đặt trước khi bật lại container. Phép so này \emph{không ngang hàng}: MySQL trong đề tài chạy một node, không có bản sao, và có thể triển khai cùng InnoDB Cluster. Cả hai hệ đều không mất giao dịch đã xác nhận. Vì vậy demo cho thấy khác biệt về tính sẵn sàng khi mất một tiến trình, không cho thấy khác biệt về độ bền.

Demo có các giới hạn sau. Ba TiKV nằm trên một máy và một ổ đĩa, nên đây là lỗi của một tiến trình, không phải lỗi của một máy; PD và TiDB server chỉ có một node và không được tắt. Mỗi hệ chỉ có một lượt trong bảng; một lượt chạy thử trước đó của TiDB cho 8 giây ngừng và 10 giây để về 90\% mức nền, nên thời gian phục hồi dao động giữa các lần. Số giao dịch mỗi giây của demo không dùng để so hiệu năng hai hệ.

\subsection{Demo bổ sung: thêm node TiKV}

Demo này thêm \code{tikv4} vào cụm C3 đang có tải ghi. Bảng \code{shipments} có 100.000 dòng và được tách sẵn thành nhiều Region để PD có dữ liệu cần cân bằng. \code{tikv4} được thêm ở giây 30.

\begin{table}[!htp]
\centering\bangfont
\caption{Phân bố leader và Region khi thêm tikv4 (leader / Region)}\label{tab:demo-them-node}
\begin{tabular}{|C{3.0cm}|C{2.6cm}|C{2.6cm}|C{2.6cm}|C{2.6cm}|}
\hline
\dongdau \textbf{Thời điểm} & \textbf{tikv1} & \textbf{tikv2} & \textbf{tikv3} & \textbf{tikv4} \\ \hline
Trước khi thêm & 37 / 99 & 27 / 99 & 35 / 99 & --- \\ \hline
Giây 60 & 27 / 85 & 24 / 94 & 28 / 86 & 20 / 32 \\ \hline
Giây 149 & 27 / 69 & 24 / 77 & 28 / 82 & 20 / 69 \\ \hline
\end{tabular}
\end{table}

\code{tikv4} ở trạng thái Up sau 11 giây và nhận Region đầu tiên sau 12 giây kể từ lúc thêm. Đến cuối lượt đo, nó giữ 69 bản sao Region và 20 leader (Bảng~\ref{tab:demo-them-node}, Hình~\ref{fig:dashboard4}). Không có lệnh ghi nào bị lỗi và không mất giao dịch nào trong 124.628 giao dịch. Trong quá trình cân bằng, thông lượng có những giây giảm mạnh, thấp nhất 38 giao dịch mỗi giây so với mức trước khi thêm là 903, và trung bình 30 giây cuối là 785. Demo cho thấy có thể mở rộng trực tuyến; nó \emph{không} chứng minh thêm node làm tăng hiệu năng, vì node mới dùng chung máy và ổ đĩa với các node cũ và tải chỉ ghi vào một bảng nhỏ.

\hinh{tidb_dashboard_4_tikv.png}{0.9}{TiDB Dashboard sau khi thêm tikv4: 4 TiKV ở trạng thái Up}{fig:dashboard4}

\subsection{Kết luận chương}

Chương này trình bày cách cài đặt, cấu hình và vận hành cụm TiDB v8.5.5 gồm 1 PD, 1 TiDB server và 3 TiKV trên một máy 16 GB, cùng MySQL 8.4.11 để đối chứng. Các chức năng cơ bản chạy đúng trên cả hai hệ. Cấu hình C3 chịu được việc mất một tiến trình TiKV và không mất giao dịch đã xác nhận trong demo, và cho phép thêm node khi đang có tải. Các demo này chỉ minh họa cơ chế; mọi kết luận về chi phí chỉ mục trong Phần 2 dựa trên các thí nghiệm riêng.


\phanbaocao{PHẦN 2}{Phân tích chuyên sâu: chi phí chỉ mục phụ trên TiDB so với MySQL}

\section{Vấn đề, cơ chế và phương pháp}
\label{sec:phuong-phap}
% Chương mở đầu Phần 2. Phụ trách: TV1.
\phutrach{TV1 -- vấn đề và cơ chế}

\subsection{Vấn đề và câu hỏi nghiên cứu}
\label{subsec:van-de-rq}

Báo cáo này nghiên cứu chi phí của chỉ mục phụ trên hai kiến trúc lưu trữ khác nhau: MySQL InnoDB sử dụng một máy chứa cây B+ (B+Tree) và TiDB sử dụng kiến trúc SQL tính toán phân tách khỏi TiKV lưu trữ key--value theo cấu trúc RocksDB LSM-Tree. Đối tượng đo là bảng \code{shipments}, mô tả các vận đơn thương mại điện tử và bao gồm khóa đơn vị, mã vận đơn, mã người gửi, kho đích, trạng thái, tiền COD, các cột địa chỉ và hai mốc thời gian.

Trong vận hành thực tế, nhiều truy vấn đi qua các thuộc tính khác với khóa chính: tìm theo mã vận đơn, lấy danh sách đơn mới của một người gửi, đếm đơn theo kho và trạng thái hoặc lọc theo khoảng tiền COD. Mỗi truy cập thường làm phát sinh hai loại chi phí. Một là chi phí \emph{đọc}: số dòng lọc ra, số lần quay về bảng, số toán tử/dữ liệu qua coprocessor và số nút KV phải gọi. Hai là chi phí \emph{ghi}: mỗi chỉ mục phụ phải được thêm, xóa hoặc cập nhật khi dữ liệu thay đổi. Đồng thời, các cặp key--value của TiDB còn phải đi qua MVCC, giao dịch hai pha và nhân bản Raft.

Bốn câu hỏi nghiên cứu xác định phạm vi phân tích và các phép đo tương ứng:
\begin{enumerate}
  \item \textbf{RQ1 -- Chỉ mục covering có làm giảm chi phí truy vấn?} So sánh truy vấn chỉ dùng khóa và cột phụ trong chỉ mục với truy vấn quay về bảng.
  \item \textbf{RQ2 -- Khi nào một truy vấn nên dùng chỉ mục thay vì quét bảng?} So sánh ba chế độ truy cập \code{FORCE INDEX}, \code{IGNORE INDEX} và optimizer tự chọn theo độ chọn lọc.
  \item \textbf{RQ3 -- Khóa chính ảnh hưởng thế nào đến chỉ mục phụ?} So sánh clustered key và nonclustered key, số lần gọi TiKV và dung lượng mục chỉ mục.
  \item \textbf{RQ4 -- Việc tăng số chỉ mục phụ làm giảm thông lượng ghi ra sao?} Đo tải chèn và cập nhật theo bộ 0, 1, 3 và 5 chỉ mục, đồng thời tách chi phí giao dịch và nhân bản.
\end{enumerate}

Các câu hỏi trên không được đặt thành giả định rằng TiDB luôn nhanh hơn hoặc chậm hơn MySQL. Mục đích là xác định \emph{điều kiện} mà mỗi cơ chế mang lại lợi ích đọc, mức dung lượng và chi phí ghi đi kèm, cũng như giới hạn của phép so sánh trên một máy lab.

\subsection{Cơ chế chỉ mục phụ}
\label{subsec:co-che-index}

\subsubsection{Lưu trữ và truy cập chỉ mục phụ trong MySQL InnoDB}

Trong InnoDB, bảng có khóa chính clustered được tổ chức thành một B+Tree. Lá của cây chứa dữ liệu toàn hàng, còn mỗi chỉ mục phụ là một B+Tree riêng. Lá chỉ mục lưu giá trị khóa, khóa chính của dòng, header bản ghi và thư mục trang. Vì một truy vấn quét chỉ mục phụ có thể cần lấy các cột nằm ngoài leaf entry, InnoDB phải dùng khóa chính để quay về clustered index.

Chẳng hạn, truy vấn \code{SELECT * FROM shipments WHERE sender\_id=?} đọc khối khóa và \code{id} ở lá chỉ mục rồi tra từng dòng trong cây dữ liệu. Truy vấn \code{SELECT sender\_id, created\_at, status, cod\_amount FROM shipments WHERE sender\_id=?} có thể trở thành \emph{covering} khi tất cả cột trả về nằm trong \code{idx\_sender\_cov}. Trường hợp này không cần quay về bảng.

Khi chèn khóa ngẫu nhiên, InnoDB có thể tách trang B+Tree. Khi khóa tăng tuần tự, các trang lá được điền gần dọc thứ tự. Khi cập nhật một cột nằm trong chỉ mục phụ, engine phải thay đổi leaf entry tương ứng; đồng thời lệnh ghi tạo redo log, có thể tạo undo record cho MVCC và chịu cơ chế group commit theo cấu hình giao dịch. Như vậy, chỉ số \emph{không còn chứa dữ liệu tại lá} hay \emph{ít ghi hơn} là kết quả của từng workload, không phải tính chất chung cho mọi truy vấn.

\subsubsection{Ánh xạ chỉ mục phụ và đường truy cập trong TiDB}

TiDB ánh xạ bảng quan hệ sang key--value. Với bảng clustered, dòng có dạng \code{t{table\_id}\_r{id}}, còn chỉ mục phụ có key chứa \code{table\_id}, \code{index\_id}, giá trị cột khóa và handle của dòng. Với chỉ mục \code{idx\_sender\_created(sender\_id, created\_at)}, key không phủ ngoài khóa còn lưu trong bảng là \code{id}; với \code{idx\_sender\_cov(sender\_id, created\_at, status, cod\_amount)}, các cột \code{status} và \code{cod\_amount} được thêm vào khóa vì cùng nằm trong leaf index entry của TiKV.

Tầng TiDB lựa chọn ba đường truy cập cơ bản:
\begin{itemize}
  \item \code{Point\_Get}: tra thẳng một khóa duy nhất hoặc khóa chính mà không cần coprocessor nhiều pha.
  \item \code{IndexReader}: đọc cột cần thiết trực tiếp từ Region chứa chỉ mục; đây là đường covering khi cột truy vấn nằm trong entry index.
  \item \code{IndexLookUp}: đọc index để lấy các handle, gom handle theo lô rồi gửi yêu cầu tới Region chứa dòng dữ liệu.
\end{itemize}

Trong \code{IndexLookUp}, index task và table task có thể nằm trên TiKV khác nhau. Bản ghi thực thi ghi riêng \code{cop\_task}, \code{rpc\_num}, \code{total\_process\_keys}, \code{tot\_proc}, \code{tot\_wait} và tổng thời gian RPC. Các đại lượng này cho phép tách chi phí quay về bảng khỏi chi phí xử lý mỗi dòng. Tuy nhiên, phần \code{rest\_ms} trong mô hình không thể coi là độ trễ mạng vật lý vì cả cụm đều chạy trong các container trên một máy.

\subsubsection{Ghi chỉ mục phụ trong TiDB}

Trong TiDB, một giao dịch nạp dòng và một giao dịch cập nhật dòng đều làm thay đổi nhiều cặp key--value. Với \code{INSERT}, một dòng và mỗi mục chỉ mục phụ phải được \code{prewrite}; sau đó giao dịch gọi TSO, ghi khóa \code{commit} và chờ xác nhận commit. Với \code{UPDATE}, phần key--value của cột không nằm trong chỉ mục được ghi mới, còn mỗi cột nằm trong chỉ mục tạo ra một cặp \emph{xóa mục cũ -- ghi mục mới} cùng một phiên bản MVCC mới.

Mỗi cặp key--value sau đó được lưu trong RocksDB, đi qua WAL/MemTable, flush thành SST và compaction LSM-Tree. Một Region chỉ commit sau khi đa số peer ghi và áp dụng Raft log. Với C3, quorum thực tế là 2 trong 3 peer; với C1, Region chỉ có một peer nên không cần chờ peer thứ hai. Vì các bước trên lặp lại theo số cột/key bị ảnh hưởng, số chỉ mục phụ không chỉ làm tăng dung lượng mà còn làm tăng số key phải giao dịch, số Raft log và công việc compaction.

Tuy nhiên, ưu thế kiến trúc LSM-Tree không phải áp dụng trực tiếp cho mọi workload. Ghi luôn tuần tự không loại bỏ chi phí giao dịch phân tán, chi phí đồng thuận hoặc chi phí compaction. RQ4 vì vậy kiểm chứng hai giả thuyết trái chiều: TiDB có thể có write penalty nhỏ hơn MySQL nhờ ghi LSM, hoặc có thể thua MySQL do overhead phân tán và nhân bản.

\subsubsection{Covering index và đánh đổi}

Một covering index là chỉ mục chứa đủ cột để thực thi truy vấn mà không cần quay về bảng. Tác dụng phụ thuộc vào ba yếu tố: số dòng cần xử lý, tỷ lệ truy vấn cần cột ở ngoài chỉ mục và chi phí truy cập thêm. Trong RQ1, \code{idx\_sender\_cov(sender\_id, created\_at, status, cod\_amount)} vừa là ví dụ covering cho truy vấn lấy 20 đơn mới nhất, vừa là một cấu trúc lớn hơn \code{idx\_sender\_created} ở mục B.1. Covering không miễn phí:
\begin{itemize}
  \item Khi truy vấn tra ít dòng, nó có thể loại bỏ hai pha \code{IndexLookUp} hoặc lượt quay về bảng và làm giảm độ trễ.
  \item Khi truy vấn đọc số dòng lớn, kích thước entry lớn hơn làm truy xuất chỉ mục truyền nhiều dữ liệu hơn và có thể không còn lợi thế.
  \item Khi \code{status} hoặc \code{cod\_amount} được cập nhật, chỉ mục covering phải ghi/xóa thêm mục; RQ4 đo ảnh hưởng tương ứng.
  \item Chỉ mục có nhiều cột làm tăng dung lượng, dung lượng tạo thêm và thời gian duy trì. Số liệu thực tế của hai hệ được đối chiếu ở Chương~\ref{sec:dung-luong}.
\end{itemize}

\begin{table}[!htp]
\centering\bangfont
\caption{So sánh cơ chế chỉ mục phụ của hai hệ}\label{tab:so-co-che}
\begin{tabular}{|L{3.4cm}|L{5.2cm}|L{5.8cm}|}
\hline
\dongdau \textbf{Khía cạnh} & \textbf{MySQL InnoDB} & \textbf{TiDB/TiKV} \\ \hline
Cấu trúc chỉ mục & B+Tree riêng, lá lưu khóa, khóa chính và dữ liệu khóa & Cặp key--value trong Region/TiKV, key chứa index ID và giá trị khóa, handle hoặc dữ liệu nhỏ \\ \hline
Đọc không covering & Quét leaf index rồi quay về clustered B+Tree bằng khóa chính & \code{IndexLookUp}: đọc index rồi gom handle và gửi table task đến Region dòng dữ liệu \\ \hline
Đọc covering & Trả kết quả trực tiếp từ leaf B+Tree nếu đủ cột & \code{IndexReader}: chỉ đọc leaf index entry qua coprocessor \\ \hline
Thêm một dòng & Thêm leaf entry vào mỗi chỉ mục; khóa tăng tuần tự hoặc khóa ngẫu nhiên ảnh hưởng tách trang & Tạo cặp key--value cho dòng và mỗi chỉ mục; đi qua prewrite, commit, MVCC và Raft \\ \hline
Cập nhật cột trong chỉ mục & Cập nhật leaf entry tương ứng, redo log và các cơ chế giao dịch của InnoDB & Xóa entry cũ và ghi entry mới cùng MVCC, sau đó ghi LSM/Raft theo số Region liên quan \\ \hline
Khóa chính dài & Khóa chính được lưu trong leaf của mọi chỉ mục phụ & Handle/key clustered hoặc \code{\_tidb\_rowid} nonclustered tăng kích thước/dải gọi KV \\ \hline
Chi phí đo chính & Thời gian truy vấn, Handler reads, dung lượng trang và byte ghi & Thời gian truy vấn, cop task/RPC/key, Region size và byte ghi đĩa \\ \hline
\end{tabular}
\end{table}

% Phương pháp triển khai và kịch bản. Phụ trách: TV2.
\phutrach{TV2 -- phương pháp triển khai và kịch bản}

\subsection{Phương pháp triển khai thực nghiệm}
\ghichu{Lệnh kiểm tra và \code{team\_hash}; sinh dữ liệu seed 20261116; chuyển hệ, nạp, checksum; kiểm tra bản sao; chờ đĩa lắng; chạy tiếp sau khi bị ngắt; chỉ so hai hệ trên cùng máy, giữa các máy chỉ so tỉ lệ. Cấu hình máy và hồ sơ XS ở Chương~\ref{sec:cai-dat}.}

\begin{table}[!htp]
\centering\bangfont
\caption{Đối chiếu cấu hình và dữ liệu giữa ba máy}\label{tab:checksum}
\begin{tabular}{|C{1.4cm}|L{3.6cm}|L{4.6cm}|L{4.6cm}|}
\hline
\dongdau \textbf{Máy} & \textbf{team\_hash} & \textbf{sum\_cod, 1 triệu dòng} & \textbf{sum\_cod, 3 triệu dòng} \\ \hline
M1 & \chuaco & \chuaco & \chuaco \\ \hline
M2 & fa4c12cdf33a & 5000968738439 & --- \\ \hline
M3 & \chuaco & \chuaco & \chuaco \\ \hline
\end{tabular}
\end{table}

\subsection{Kịch bản thử nghiệm}
Bảng~\ref{tab:kich-ban} tóm tắt các kịch bản. Thiết kế chi tiết nằm ở đầu mỗi chương kết quả.

\begin{table}[!htp]
\centering\bangnho
\caption{Tổng hợp các kịch bản thử nghiệm}\label{tab:kich-ban}
\begin{tabular}{|C{1.3cm}|L{3.4cm}|C{1.5cm}|L{2.4cm}|L{3.0cm}|L{2.4cm}|}
\hline
\dongdau \textbf{Mã} & \textbf{Biến thay đổi} & \textbf{Máy} & \textbf{Số dòng} & \textbf{Chỉ số} & \textbf{Trình bày} \\ \hline
B.1 & Số chỉ mục (0, 5) & M1 & 1 triệu, 3 triệu & Dung lượng, tài nguyên tạo chỉ mục & Chương~\ref{sec:dung-luong} \\ \hline
RQ1 & Covering, R, truy vấn & M3, M1 & 1 triệu (M3), 3 triệu (M1) & Kế hoạch, độ trễ, cop task, RPC & Chương~\ref{sec:doc} \\ \hline
RQ2 & Độ chọn lọc, chế độ ép kế hoạch & M3 & 3 triệu & Độ trễ, điểm giao nhau & Chương~\ref{sec:doc} \\ \hline
RQ4 & Số chỉ mục, tải, cấu hình bản sao & M2 & 1 triệu & TPS, p95, p99, byte ghi & Chương~\ref{sec:ghi} \\ \hline
\end{tabular}
\end{table}

\subsection{Cách xử lý số liệu}
\ghichu{Trung vị và min--max của các vòng; không bỏ lượt; lượt bất thường được nêu riêng; mỗi CSV ghi máy, hồ sơ, cấu hình bản sao và mã cấu hình.}


\clearpage
\section{Thực nghiệm sơ bộ, dữ liệu và dung lượng chỉ mục}
\label{sec:dung-luong}
% Mục B.0. Phụ trách: TV3.
\phutrach{TV3 -- mục B.0}

\subsection{Thực nghiệm sơ bộ}
\ghichu{Khoảng nửa trang. sysbench trên \code{sbtest1}, 16 đến 512 luồng. Nêu giả thuyết ban đầu không được xác nhận. Chỉ giữ số có log gốc; ghi rõ đây là số giai đoạn 1, khác cấu hình XS. Kết thúc bằng câu hỏi dẫn tới các mục sau.}

\begin{table}[!htp]
\centering\bangfont
\caption{Thông lượng sơ bộ trên sbtest1 (giai đoạn 1)}\label{tab:so-bo}
\begin{tabular}{|L{4.4cm}|C{2.2cm}|C{2.6cm}|C{2.6cm}|C{2.2cm}|}
\hline
\dongdau \textbf{Tải} & \textbf{Số luồng} & \textbf{MySQL} & \textbf{TiDB} & \textbf{Tỉ số} \\ \hline
\code{oltp\_read\_write} & \chuaco & \chuaco & \chuaco & \chuaco \\ \hline
\code{oltp\_update\_index} & \chuaco & \chuaco & \chuaco & \chuaco \\ \hline
\end{tabular}
\end{table}

% Mục B.1. Phụ trách: TV1.
% Số liệu: results/b1_storage_M2.csv (đo), results/b1_model.csv (mô hình), results/b1_summary.csv.
\phutrach{TV1 -- mục B.1}

\subsection{Dữ liệu và lược đồ}
Bảng \code{shipments} có 14 cột, sinh bằng seed 20261116. Hai hệ dùng cùng lược đồ; TiDB thêm từ khóa \code{CLUSTERED} cho khóa chính.

\begin{lstlisting}[language=SQL,caption={Lược đồ bảng shipments}]
CREATE TABLE shipments (
  id BIGINT NOT NULL AUTO_INCREMENT,
  tracking_code VARCHAR(20) NOT NULL,
  sender_id INT NOT NULL,
  origin_hub_id SMALLINT NOT NULL,
  dest_hub_id SMALLINT NOT NULL,
  status TINYINT NOT NULL,
  cod_amount INT NOT NULL DEFAULT 0,
  weight_grams INT NOT NULL,
  shipping_fee INT NOT NULL,
  recipient_phone VARCHAR(15) NOT NULL,
  recipient_address VARCHAR(150) NOT NULL,
  note VARCHAR(100),
  created_at DATETIME NOT NULL,
  updated_at DATETIME NOT NULL,
  PRIMARY KEY (id)   -- TiDB: PRIMARY KEY (id) CLUSTERED
);
\end{lstlisting}

\begin{table}[!htp]
\centering\bangfont
\caption{Các bộ chỉ mục dùng trong B.1 và RQ4}\label{tab:bo-chi-muc}
\begin{tabular}{|L{3.0cm}|L{11.8cm}|}
\hline
\dongdau \textbf{Bộ} & \textbf{Chỉ mục} \\ \hline
0 & Không có chỉ mục phụ \\ \hline
1 & \code{idx\_code} (tracking\_code), UNIQUE \\ \hline
3 & Bộ 1; \code{idx\_sender\_created} (sender\_id, created\_at); \code{idx\_hub\_status} (dest\_hub\_id, status) \\ \hline
5 & Bộ 3; \code{idx\_cod} (cod\_amount); \code{idx\_status} (status) \\ \hline
sender\_created & Chỉ \code{idx\_sender\_created} \\ \hline
sender\_cov & Chỉ \code{idx\_sender\_cov} (sender\_id, created\_at, status, cod\_amount) \\ \hline
\end{tabular}
\end{table}

Trong dữ liệu, mã vận đơn luôn dài 14 ký tự, số điện thoại 10 ký tự, địa chỉ khoảng 60 ký tự, và khoảng 40\% đơn không có ghi chú (\code{note} là NULL). Mọi chuỗi đều là ký tự ASCII, nên mỗi ký tự chiếm 1 byte trên cả hai hệ.

\subsection{Mô hình dung lượng của InnoDB}
Mô hình tính dung lượng của từng bản ghi theo định dạng DYNAMIC của InnoDB~\cite{mysql_innodb}, rồi cộng lại cho mọi dòng của file dữ liệu. Một bản ghi của cây khóa chính gồm 5 B đầu bản ghi, 1 B độ dài cho mỗi cột \code{VARCHAR} không NULL (mọi giá trị ở đây đều ngắn hơn 128 B), 1 B đánh dấu NULL, 6 B mã giao dịch và 7 B con trỏ hoàn tác ẩn, cùng các cột: số nguyên theo đúng kiểu (8, 4, 2 hoặc 1 B), \code{DATETIME} 5 B và chuỗi theo độ dài thật. Một bản ghi của chỉ mục phụ chỉ gồm 5 B đầu bản ghi, các cột khóa và khóa chính 8 B, nên dài từ 14 B (\code{idx\_status}) đến 28 B (\code{idx\_code}).

Bản ghi được xếp vào trang 16 KiB, mỗi trang mất 128 B cho phần đầu và cuối trang, cộng khoảng 2 B thư mục trang cho mỗi 6 bản ghi. Trang của cây khóa chính được lấp bằng các lệnh chèn theo \code{id} tăng dần nên chừa trống 1/16; trang của chỉ mục phụ được dựng bằng cách sắp xếp (\code{innodb\_fill\_factor=100}) nên được lấp đầy. Mô hình bỏ qua các trang không phải lá, vốn chỉ chiếm khoảng 1\%.

\subsection{Mô hình dung lượng của TiDB/TiKV}
Mô hình mã hóa từng dòng và từng mục chỉ mục đúng như TiDB ghi xuống TiKV (Bảng~\ref{tab:anh-xa-kv}), và tính kích thước logic của một bản sao, trước khi nén. Một khóa trong TiKV gồm 1 B tiền tố, khóa của TiDB được mã hóa để so sánh được theo byte (cứ 8 B dữ liệu thêm 1 B đánh dấu), và 8 B dấu thời gian MVCC~\cite{tikv_storage}. Giá trị không quá 255 B được lưu ngay trong column family \code{write}, kèm 1 B loại, 9 B \code{start\_ts} và 2 B tiền tố; giá trị dài hơn được chuyển sang column family \code{default}, và khóa phải lặp lại ở đó.

Giá trị của một dòng dùng định dạng dòng thế hệ 2: 6 B đầu, 1 B mã cột cho mỗi cột, 2 B vị trí cho mỗi cột không NULL, số nguyên được thu gọn còn 1, 2, 4 hoặc 8 B, \code{DATETIME} 8 B và chuỗi giữ nguyên. Trong khóa của chỉ mục, mọi số nguyên và \code{DATETIME} đều chiếm 9 B, bất kể kiểu khai báo. Mục của chỉ mục thường chứa handle của dòng trong khóa và có giá trị 1 B; mục của chỉ mục UNIQUE được giả định lưu handle 8 B trong giá trị.

\subsection{Phương pháp đo}
Số đo lấy trên máy M2, hồ sơ XS, cấu hình C3. Mỗi quy mô dùng một bảng mới trong database riêng: nạp dữ liệu, kiểm tra số dòng và tổng \code{cod\_amount}, rồi tạo lần lượt từng chỉ mục của bộ 5 và cuối cùng là \code{idx\_sender\_cov}. Dung lượng của mỗi chỉ mục được đọc ngay sau khi tạo.

Trên MySQL, dung lượng lấy từ \code{mysql.innodb\_index\_stats} sau \code{ANALYZE TABLE}, theo hai con số: số trang đã cấp phát cho chỉ mục (dung lượng thật trong file dữ liệu) và số trang lá đang dùng (để so với mô hình).

Trên TiDB, nhóm không dùng thống kê của PD hay \code{TABLE\_STORAGE\_STATS}, vì khi chạy thử các số này chỉ được cập nhật khi Region có thay đổi nên bị cũ. Thay vào đó, dòng và từng chỉ mục được tách thành các Region riêng; chương trình kiểm tra mọi Region có dữ liệu chỉ chứa khóa của đúng đối tượng cần đo, rồi đọc kích thước từng Region trên store \code{tikv1} bằng \code{tikv-ctl size}. Con số này là kích thước logic của các cặp khóa--giá trị trong file SST, cho một bản sao, trước khi nén. Vì dòng được nạp bằng giao dịch thông thường, một phần có thể còn nằm trong bộ nhớ, nên \code{tikv1} được compaction trước khi đo; sau compaction, mỗi dòng chỉ còn một phiên bản MVCC. Chỉ mục được TiDB dựng bằng cách nạp thẳng file SST nên không gặp vấn đề này. Dung lượng toàn cụm C3 xấp xỉ gấp ba con số một bản sao; dung lượng sau nén trên đĩa không được đo.

\subsection{Đối chiếu mô hình và số đo}
\begin{table}[!htp]
\centering\bangnho
\caption{Dung lượng theo mô hình và theo số đo (MiB, một bản sao)}\label{tab:dung-luong}
\begin{tabular}{|L{2.3cm}|C{1.1cm}|C{1.5cm}|C{1.9cm}|C{1.6cm}|C{1.5cm}|C{1.5cm}|C{1.5cm}|}
\hline
\dongdau \textbf{Quy mô, bộ chỉ mục} & \textbf{Thô} & \textbf{MySQL mô hình} & \textbf{MySQL đo: cấp phát (trang lá)} & \textbf{MySQL trang lá / mô hình} & \textbf{TiDB mô hình} & \textbf{TiDB đo} & \textbf{TiDB đo / mô hình} \\ \hline
1 triệu, 0 & 136,3 & 170,0 & 171,7 (171,3) & 1,007 & 216,5 & 216,2 & 0,998 \\ \hline
1 triệu, 5 & 136,3 & 264,9 & 283,5 (267,1) & 1,008 & 559,9 & 582,4 & 1,040 \\ \hline
3 triệu, 0 & 408,8 & 510,1 & 516,0 (513,8) & 1,007 & 649,6 & 648,5 & 0,998 \\ \hline
3 triệu, 5 & 408,8 & 794,6 & 847,1 (801,1) & 1,008 & 1.679,6 & 1.747,1 & 1,040 \\ \hline
10 triệu, 5 (mô hình) & 1.362,8 & 2.648,6 & --- & --- & 5.598,6 & --- & --- \\ \hline
\end{tabular}
\par\vspace{3pt}{\noindent\small\itshape ``Thô'' là tổng kích thước các cột theo kiểu của MySQL, không kể phần phụ. Máy M2, cấu hình C3.\par}
\end{table}

Với MySQL, mô hình gần như trùng với số trang lá đo được: lệch 0,7\% với phần dữ liệu và 0,7--1,2\% với từng chỉ mục, ở cả hai quy mô (Bảng~\ref{tab:dung-luong}). Dung lượng cấp phát thì lớn hơn: chỉ mục phụ chiếm thêm 16--20\% so với mô hình, vì InnoDB cấp trang theo từng nhóm và giữ sẵn một phần chưa dùng. Ở bộ 5 chỉ mục, phần chênh này làm tổng dung lượng cấp phát lớn hơn mô hình khoảng 7\%.

Với TiDB, mô hình khớp số đo ở phần dữ liệu (lệch 0,2\%) và ở năm chỉ mục thường (đúng tới từng byte cho mỗi mục). Riêng \code{idx\_code}, chỉ mục UNIQUE duy nhất, đo được 98,0 B mỗi mục trong khi mô hình cho 74,0 B: giá trị của mục dài khoảng 32 B chứ không phải 8 B như giả định. Nhóm chưa giải mã giá trị này nên không biết 24 B dư ra gồm những gì; đây là phần chênh duy nhất của mô hình và là lý do tổng bộ 5 của TiDB lớn hơn mô hình 4\%.

Bảng~\ref{tab:muc-chi-muc} so kích thước một mục chỉ mục trên hai hệ. Mục của TiDB lớn hơn MySQL 3,4 đến 4,8 lần. Phần lớn chênh lệch là chi phí cố định của mỗi cặp khóa--giá trị: tiền tố bảng và chỉ mục, byte đánh dấu của mã hóa so sánh được, 8 B dấu thời gian và 13 B phần đầu giá trị. Vì mọi số nguyên trong khóa đều chiếm 9 B, \code{idx\_hub\_status} trên hai cột \code{SMALLINT} và \code{TINYINT} có mục lớn đúng bằng \code{idx\_sender\_created}; trên MySQL, mục của nó nhỏ hơn 6 B.

\begin{table}[!htp]
\centering\bangfont
\caption{Kích thước một mục chỉ mục (B)}\label{tab:muc-chi-muc}
\begin{tabular}{|L{4.2cm}|C{2.6cm}|C{2.6cm}|C{2.6cm}|C{2.4cm}|}
\hline
\dongdau \textbf{Chỉ mục} & \textbf{MySQL, bản ghi (mô hình)} & \textbf{TiDB mô hình} & \textbf{TiDB đo} & \textbf{TiDB đo / MySQL} \\ \hline
\code{idx\_code} (UNIQUE) & 28 & 74 & 98 & 3,5 \\ \hline
\code{idx\_sender\_created} & 22 & 76 & 76 & 3,5 \\ \hline
\code{idx\_hub\_status} & 16 & 76 & 76 & 4,8 \\ \hline
\code{idx\_cod} & 17 & 67 & 67 & 3,9 \\ \hline
\code{idx\_status} & 14 & 67 & 67 & 4,8 \\ \hline
\code{idx\_sender\_cov} & 27 & 103 & 103 & 3,8 \\ \hline
\end{tabular}
\par\vspace{3pt}{\noindent\small\itshape Bản ghi MySQL chưa tính phần thư mục trang và phần trống của trang. Số đo TiDB là dung lượng chỉ mục chia số dòng, giống nhau ở 1 và 3 triệu dòng.\par}
\end{table}

Ở quy mô 10 triệu dòng, mô hình cho bộ 5 chỉ mục khoảng 2,6 GiB trên MySQL và 5,5 GiB cho một bản sao trên TiDB, tức khoảng 16,4 GiB logic cho cả cụm C3 trước khi nén. Đây chỉ là số ngoại suy từ mô hình, không phải số đo.

\subsection{Chi phí tạo chỉ mục}
\begin{table}[!htp]
\centering\bangnho
\caption{Chi phí tạo từng chỉ mục ở 3 triệu dòng}\label{tab:tao-chi-muc}
\begin{tabular}{|L{3.4cm}|C{1.5cm}|C{1.6cm}|C{1.8cm}|C{1.5cm}|C{1.6cm}|C{1.8cm}|}
\hline
\dongdau \textbf{Chỉ mục} & \textbf{MySQL (s)} & \textbf{MySQL (CPU)} & \textbf{MySQL ghi (MiB)} & \textbf{TiDB (s)} & \textbf{TiDB (CPU)} & \textbf{TiDB ghi (MiB)} \\ \hline
\code{idx\_code} & 4,8 & 1,51 & 962 & 27,0 & 1,23 & 618 \\ \hline
\code{idx\_sender\_created} & 2,9 & 1,20 & 731 & 16,0 & 1,69 & 148 \\ \hline
\code{idx\_hub\_status} & 2,6 & 1,36 & 501 & 14,8 & 1,72 & 84 \\ \hline
\code{idx\_cod} & 2,6 & 1,22 & 544 & 12,3 & 1,69 & 123 \\ \hline
\code{idx\_status} & 2,3 & 1,28 & 429 & 11,8 & 1,64 & 40 \\ \hline
\code{idx\_sender\_cov} & 4,9 & 0,90 & 918 & 17,8 & 1,57 & 295 \\ \hline
\end{tabular}
\par\vspace{3pt}{\noindent\small\itshape CPU: số CPU dùng trung bình trong lúc tạo (ngân sách mỗi hệ là 4). Ghi: byte ghi đĩa của các container cơ sở dữ liệu, đo một lần cho mỗi chỉ mục.\par}
\end{table}

Ở 3 triệu dòng, MySQL tạo mỗi chỉ mục trong 2,3--4,9 giây, còn TiDB cần 11,8--27,0 giây, tức chậm hơn 3,6 đến 5,7 lần (Bảng~\ref{tab:tao-chi-muc}). TiDB dùng 1,2--1,7 CPU trên ngân sách 4, nên tốc độ tạo chỉ mục không bị giới hạn bởi tổng CPU của cụm. Lượng byte ghi đĩa của MySQL lớn hơn nhiều lần dung lượng của chỉ mục; nhóm chưa tách phần này thành file sắp xếp tạm, redo log và trang dữ liệu. Byte ghi của TiDB là byte đã nén trên đĩa của cả ba bản sao, nên nhỏ hơn dung lượng logic trong Bảng~\ref{tab:dung-luong}. Riêng \code{idx\_code} ghi nhiều hơn hẳn các chỉ mục khác của TiDB, có thể vì đây là chỉ mục được tạo đầu tiên trên bảng; nhóm chưa kiểm tra nguyên nhân. Mỗi chỉ mục chỉ được tạo một lần, nên các con số này chỉ cho biết độ lớn.

\subsection{Dung lượng của covering index}
\begin{table}[!htp]
\centering\bangfont
\caption{Dung lượng chỉ mục có và không covering (MiB, một bản sao)}\label{tab:dl-covering}
\begin{tabular}{|L{4.2cm}|C{2.0cm}|C{3.0cm}|C{3.0cm}|C{2.2cm}|}
\hline
\dongdau \textbf{Chỉ mục} & \textbf{Số dòng} & \textbf{MySQL: cấp phát (trang lá)} & \textbf{TiDB} & \textbf{Byte mỗi mục, MySQL / TiDB} \\ \hline
\code{idx\_sender\_created} & 1 triệu & 25,6 (21,7) & 72,5 & 22 / 76 \\ \hline
\code{idx\_sender\_cov} & 1 triệu & 30,6 (26,5) & 98,2 & 27 / 103 \\ \hline
\code{idx\_sender\_created} & 3 triệu & 74,6 (64,9) & 217,4 & 22 / 76 \\ \hline
\code{idx\_sender\_cov} & 3 triệu & 91,7 (79,5) & 294,7 & 27 / 103 \\ \hline
\end{tabular}
\end{table}

Thêm hai cột \code{status} và \code{cod\_amount} để chỉ mục trở thành covering làm dung lượng tăng khoảng 22\% trên MySQL (theo trang lá) và 36\% trên TiDB, giống nhau ở cả hai quy mô (Bảng~\ref{tab:dl-covering}). Về tuyệt đối, ở 3 triệu dòng, phần thêm là 14,6 MiB trên MySQL và 77,3 MiB cho mỗi bản sao trên TiDB. Mỗi cột số nguyên thêm vào khóa tốn 4 hoặc 1 B trên MySQL, nhưng 9 B trên TiDB, cộng với phần mở rộng của mã hóa so sánh được.

\subsection{Kết luận chương}
Mô hình dựa trên định dạng lưu trữ dự đoán được dung lượng của cả hai hệ khá sát. Với MySQL, mô hình trùng với số trang lá trong khoảng 1,2\%, nhưng dung lượng cấp phát thật của chỉ mục phụ lớn hơn 16--20\%. Với TiDB, mô hình khớp tới từng byte cho dữ liệu và chỉ mục thường, và chỉ sai ở chỉ mục UNIQUE, nơi giá trị của mục dài hơn giả định 24 B mà nhóm chưa giải thích được. Trên cùng dữ liệu, mỗi mục chỉ mục của TiDB lớn hơn MySQL 3,4 đến 4,8 lần, và cả cụm C3 còn lưu ba bản sao, nên chi phí dung lượng của mỗi chỉ mục phụ trên TiDB cao hơn rõ rệt. Các số này là kích thước logic trước khi nén; dung lượng sau nén trên đĩa của TiDB chưa được đo.


\clearpage
\section{Chi phí đọc qua chỉ mục phụ}
\label{sec:doc}
% Mục B.2. Phụ trách: TV3. Số RQ1 ở 3 triệu dòng do TV1 gửi (results/rq1_M1*.csv).
\phutrach{TV3 -- mục B.2}

\subsection{RQ1: thiết kế thí nghiệm}

\subsubsection{Bộ truy vấn}
\begin{table}[!htp]
\centering\bangnho
\caption{Bộ truy vấn vận đơn, thiết kế không covering và covering}\label{tab:truy-van}
\begin{tabular}{|C{1.1cm}|L{3.4cm}|L{4.1cm}|L{4.1cm}|C{1.6cm}|}
\hline
\dongdau \textbf{Mã} & \textbf{Nghiệp vụ} & \textbf{Không covering} & \textbf{B: covering} & \textbf{R} \\ \hline
Q1 & Tra một vận đơn & A: UNIQUE (tracking\_code); A2: cùng cột, không unique & (tracking\_code, status, updated\_at) & 1 \\ \hline
Q2 & 20 đơn mới nhất của một shop & A: (sender\_id, created\_at) & (sender\_id, created\_at, status, cod\_amount) & 20 \\ \hline
Q2R & Như Q2, bỏ LIMIT & Như Q2 & Như Q2 & 10 đến 10.000 \\ \hline
Q3 & Đếm đơn theo trạng thái tại một kho & A: (dest\_hub\_id) & (dest\_hub\_id, status) & khoảng N/200 \\ \hline
\end{tabular}
\end{table}

\subsubsection{Tách ba nguồn chênh lệch}
Một chênh lệch độ trễ giữa hai thiết kế chỉ mục có thể đến từ ba nguồn khác nhau, nên mỗi nguồn được đo riêng thay vì gán tất cả cho mạng.
\begin{itemize}
  \item \textbf{Độ chọn lọc} là số dòng $R$ phải đọc. Hai thiết kế trong một cặp so sánh luôn trả về cùng các dòng; Q2R chọn mốc thời gian để $R$ xấp xỉ 10, 100, 1.000 và 10.000.
  \item \textbf{Đường truy cập} là việc có phải quay về bảng hay không. Chỉ mục được ép bằng \code{FORCE INDEX} và kế hoạch được kiểm tra bằng \code{EXPLAIN} trước khi đo; nếu khác dự kiến thì dừng. Toán tử thực tế của mọi lượt nằm ở Bảng~\ref{tab:ke-hoach}.
  \item \textbf{Chi phí gọi KV} là phần cố định của mỗi pha truy cập, đo bằng số cop task, số RPC, thời gian xử lý và thời gian chờ trong \code{EXPLAIN ANALYZE} của TiDB. Phần còn lại là hiệu \code{rpc\_time} trừ xử lý trừ chờ; đó chỉ là xấp xỉ, gồm mã hóa, gRPC và lập lịch, không phải độ trễ mạng vì mọi node nằm trên một máy.
\end{itemize}
Hai đường chuẩn là \code{SELECT 1} (MySQL 0,362 ms, TiDB 0,571 ms) và tra theo khóa chính (MySQL 0,406 ms, TiDB 0,950 ms), đo trên M2 ở 3 triệu dòng. Bộ nhớ đệm kết quả của TiDB bị tắt trong lúc đo. Với Q1, thiết kế A là UNIQUE nên TiDB dùng \code{Point\_Get} thay vì coprocessor; vì vậy hiệu ứng covering của Q1 lấy từ A2 (cùng khóa nhưng không UNIQUE) so với B, còn A so với B trộn hai cơ chế khác nhau.

\subsection{RQ1: kế hoạch thực thi}
\begin{table}[!htp]
\centering\bangnho
\caption{Toán tử truy cập quan sát được}\label{tab:ke-hoach}
\begin{tabular}{|C{1.6cm}|C{2.4cm}|L{5.0cm}|L{5.0cm}|}
\hline
\dongdau \textbf{Truy vấn} & \textbf{Thiết kế} & \textbf{MySQL} & \textbf{TiDB C3} \\ \hline
Q1 & A / A2 / B & A: \code{IndexAndRow(const)}; A2: \code{IndexAndRow(ref)}; B: \code{CoveringIndex} & A: \code{Point\_Get}; A2: \code{IndexLookUp}; B: \code{IndexReader} \\ \hline
Q2, Q2R & A / B & A: \code{IndexAndRow(range)}; B: \code{CoveringIndex} & A: \code{IndexLookUp}; B: \code{IndexReader} \\ \hline
Q3 & A / B & A: \code{IndexAndRow(ref)}; B: \code{CoveringIndex} & A: \code{IndexLookUp}; B: \code{IndexReader} \\ \hline
\end{tabular}
\end{table}
Hình~\ref{fig:explain-tidb} và~\ref{fig:explain-mysql} chụp \code{EXPLAIN ANALYZE} của Q2 trên M2 (3 triệu dòng, một truy vấn mẫu, không phải số đo): trên TiDB, thiết kế A là \code{IndexLookUp} hai cop task (đọc chỉ mục rồi đọc dòng), thiết kế B là \code{IndexReader} một cop task; trên MySQL, thiết kế B là \code{Covering index range scan}.

\hinh{minh_chung_explain_tidb.png}{1.0}{Kế hoạch thực thi Q2 trên TiDB, thiết kế A (idx\_sender\_created) và B (idx\_sender\_cov); ảnh chụp terminal trên M2, một truy vấn mẫu}{fig:explain-tidb}

\hinh{minh_chung_explain_mysql.png}{1.0}{Kế hoạch thực thi Q2 trên MySQL, thiết kế A và B; ảnh chụp terminal trên M2, một truy vấn mẫu}{fig:explain-mysql}

Trong tệp \code{EXPLAIN ANALYZE} của lượt đo trên M3 (\path{results/raw}, lần lặp 1, 1 triệu dòng; truy vấn Q3 với \code{dest\_hub\_id=149}), một truy vấn mẫu cho thấy thêm:
\begin{itemize}
  \item \textbf{Q3, TiDB:} cả hai thiết kế đều quét chỉ mục 5\,003 khóa, nhưng thiết kế A còn phải đọc bảng qua 7 cop task (\code{tot\_proc} 253,3 ms, tổng \code{rpc\_info} 283 ms), còn thiết kế B gộp phép đếm xuống tầng chỉ mục với một cop task (\code{tot\_proc} 1,44 ms). Thời gian truy vấn mẫu là 88,3 ms so với 2,27 ms. Trên MySQL là 7,15 ms so với 0,608 ms: thiết kế A phải tạo bảng tạm để gom nhóm theo \code{status}, còn B đọc \code{status} thẳng từ chỉ mục và dùng \code{Group aggregate}.
\end{itemize}

\subsection{RQ1: lợi ích của covering}
Ở 3 triệu dòng trên M2 (số chính), lợi ích covering phụ thuộc vào $R$ và khác nhau giữa hai hệ.
\begin{itemize}
  \item \textbf{$R$ nhỏ} (Q1, Q2 và Q2R với $R = 10$): TiDB lợi hơn rõ, với tỉ số 1,48, 2,24 và 2,01; MySQL gần như không đổi, 0,99, 1,12 và 1,08.
  \item \textbf{$R = 100$}: 2,48 ở TiDB so với 1,56 ở MySQL.
  \item \textbf{$R \geq 1.000$}: MySQL lợi hơn, 3,56 so với 1,95 ở $R = 1.000$ và 5,91 so với 3,31 ở $R = 10.000$.
  \item \textbf{Q3} ($R = 15.000$, đếm theo trạng thái): hai hệ gần như bằng nhau, 16,78 ở MySQL và 16,86 ở TiDB.
  \item \textbf{Q1 với chỉ mục UNIQUE} (thiết kế A, \code{Point\_Get} trên TiDB): 1,608 ms, gần bằng thiết kế B (1,626 ms), nên tra một khóa duy nhất không cần covering.
\end{itemize}

\begin{table}[!htp]
\centering\bangnho
\caption[Độ trễ p50 và lợi ích covering trong từng hệ, 3 triệu dòng (M2)]{Độ trễ p50 và lợi ích covering $t(A)/t(B)$ trong từng hệ, 3 triệu dòng (M2)}\label{tab:covering}
\begin{tabular}{|C{1.5cm}|C{1.4cm}|C{1.8cm}|C{1.8cm}|C{1.8cm}|C{1.8cm}|C{1.6cm}|C{1.6cm}|}
\hline
\dongdau \textbf{Truy vấn} & \textbf{R} & \textbf{MySQL A (ms)} & \textbf{MySQL B (ms)} & \textbf{TiDB A (ms)} & \textbf{TiDB B (ms)} & \textbf{Lợi ích MySQL} & \textbf{Lợi ích TiDB} \\ \hline
Q1 (A2) & 1 & 0,429 & 0,434 & 2,408 & 1,626 & 0,99 & 1,48 \\ \hline
Q2 & 20 & 0,521 & 0,464 & 3,942 & 1,761 & 1,12 & 2,24 \\ \hline
Q2R & 10 & 0,498 & 0,462 & 3,718 & 1,847 & 1,08 & 2,01 \\ \hline
Q2R & 100 & 0,833 & 0,535 & 4,726 & 1,908 & 1,56 & 2,48 \\ \hline
Q2R & 1.000 & 3,309 & 0,929 & 7,791 & 3,987 & 3,56 & 1,95 \\ \hline
Q2R & 10.000 & 22,353 & 3,785 & 44,301 & 13,394 & 5,91 & 3,31 \\ \hline
Q3 & 15.000 & 28,266 & 1,685 & 111,876 & 6,634 & 16,78 & 16,86 \\ \hline
\end{tabular}
\par\vspace{3pt}{\noindent\small\itshape Nguồn: results/rq1\_M2\_summary.csv (số chính). Chỉ so lợi ích trong từng hệ.\par}
\end{table}

\hinh{rq1_M2_covering.png}{0.95}{Lợi ích covering theo truy vấn và $R$, 3 triệu dòng (M2)}{fig:covering}
Cùng phép đo ở 1 triệu dòng trên M3 (số liệu trong \path{results/rq1_M3_summary.csv}) cho cùng dạng với M2 ở các điểm có $R$ vừa và nhỏ: TiDB lợi hơn ở $R$ nhỏ (1,77 ở Q2 và 1,72 ở Q2R với $R = 10$, so với 1,10 và 1,07 của MySQL), còn MySQL lợi hơn từ $R = 1.000$ (2,41 so với 1,92). Ở $R = 10.000$ hai hệ gần nhau hơn trên M3 (4,03 ở MySQL, 3,80 ở TiDB) so với M2 (5,91 và 3,31). Q1 với chỉ mục không UNIQUE (A2) lợi 1,37 ở TiDB và 1,00 ở MySQL; Q1 với chỉ mục UNIQUE (\code{Point\_Get} trên TiDB) cho 0,94, tức không cần covering.

\textbf{Q3 trên TiDB không đáng tin để so.} Ở thiết kế A, p50 của 5 lượt là 99,4; 90,6; 102,8; 31,6 và 31,6 ms, và ngay trong một lượt, 5 truy vấn mẫu cho từ 15,4 đến 96,1 ms với cùng một kế hoạch; số lần đọc block ngoài bộ nhớ đệm bằng 0 ở cả năm lượt. Q2R với $R = 10.000$ ở thiết kế A cũng giảm theo cùng chiều (từ 91,1 xuống 42,8 ms). Nhóm không xác định được nguyên nhân. Tỉ số 24,6 phụ thuộc vào ba lượt chậm; nếu dùng hai lượt nhanh thì chỉ khoảng 8,6, gần bằng 8,6 của MySQL, nên báo cáo không kết luận hệ nào lợi hơn ở Q3 trên M3.

\subsection{RQ1: chi phí cố định, chi phí mỗi dòng và chi phí gọi KV}
\begin{table}[!htp]
\centering\bangfont
\caption{Mô hình $t = a + bR$ trên Q2R}\label{tab:mo-hinh}
\begin{tabular}{|C{2.4cm}|C{1.8cm}|C{1.8cm}|C{1.9cm}|C{2.2cm}|C{1.4cm}|C{2.0cm}|}
\hline
\dongdau \textbf{Máy, số dòng} & \textbf{Hệ} & \textbf{Thiết kế} & \textbf{a (ms)} & \textbf{b ($\mu$s/dòng)} & \textbf{$R^2$} & \textbf{Phần dư lớn nhất} \\ \hline
M2, 3 triệu & MySQL & A & 0,518 & 2,402 & 0,985 & 13\% \\ \hline
M2, 3 triệu & MySQL & B & 0,486 & 0,348 & 0,995 & 13\% \\ \hline
M2, 3 triệu & TiDB C3 & A & 3,918 & 3,977 & 0,999 & 14\% \\ \hline
M2, 3 triệu & TiDB C3 & B & 1,894 & 1,246 & 0,983 & 23\% \\ \hline
M1, 3 triệu & MySQL & A & 0,677 & 2,552 & 0,967 & 43\% \\ \hline
M1, 3 triệu & MySQL & B & 0,662 & 0,464 & 0,736 & 46\% \\ \hline
M1, 3 triệu & TiDB C3 & A & 7,476 & 7,668 & 0,907 & 45\% \\ \hline
M1, 3 triệu & TiDB C3 & B & 3,337 & 2,633 & 0,962 & 43\% \\ \hline
M3, 1 triệu & MySQL & A & 0,564 & 1,679 & 0,990 & 18\% \\ \hline
M3, 1 triệu & MySQL & B & 0,540 & 0,364 & 0,994 & 18\% \\ \hline
M3, 1 triệu & TiDB C3 & A & 3,095 & 4,619 & 0,818 & 46\% \\ \hline
M3, 1 triệu & TiDB C3 & B & 2,024 & 1,164 & 0,985 & 25\% \\ \hline
\end{tabular}
\par\vspace{3pt}{\noindent\small\itshape Khớp trên Q2R ($R$ từ 10 đến 10.000), trọng số $1/t$. Nguồn: results/rq1\_M2\_fit.csv, rq1\_M1\_fit.csv và rq1\_M3\_fit.csv. Không gộp các máy.\par}
\end{table}

Mô hình hai tham số giải thích dạng phụ thuộc vào $R$ này. Ở $R$ nhỏ, tỉ số tiến về $a_A / a_B$; ở $R$ lớn, tiến về $b_A / b_B$. Trên M2, $a_A / a_B$ là 1,07 ở MySQL và 2,07 ở TiDB; $b_A / b_B$ là 6,89 và 3,19. Nói cách khác, quay về bảng làm tăng chi phí \emph{cố định} của TiDB thêm 2,02 ms (MySQL chỉ thêm 0,033 ms), còn chi phí \emph{mỗi dòng} tăng thêm 2,73 $\mu$s ở TiDB và 2,05 $\mu$s ở MySQL. Chi phí mỗi dòng của đường covering trên TiDB (1,25 $\mu$s) cao hơn nhiều so với MySQL (0,35 $\mu$s), nên tỉ số tiệm cận của TiDB thấp hơn.

Các hệ số $a$, $b$ chỉ là phép khớp trong dải $R$ từ 10 đến 10.000 (độ lệch tương đối lớn nhất 23\% trên M2), không nên ngoại suy ra ngoài dải này.

\begin{table}[!htp]
\centering\bangnho
\caption{Chi phí gọi KV của TiDB theo truy vấn và thiết kế (3 triệu dòng, M2)}\label{tab:kv}
\begin{tabular}{|C{2.4cm}|C{1.3cm}|C{1.3cm}|C{1.4cm}|C{1.5cm}|C{1.7cm}|C{1.7cm}|C{2.0cm}|}
\hline
\dongdau \textbf{Truy vấn, R} & \textbf{Thiết kế} & \textbf{Cop task} & \textbf{Số RPC} & \textbf{Số key} & \textbf{Xử lý (ms)} & \textbf{Chờ (ms)} & \textbf{Phần còn lại (ms)} \\ \hline
Q1, R=1 & A2 & 2,0 & 2,0 & 2 & 0,16 & 0,11 & 0,72 \\ \hline
 & B & 1,0 & 1,0 & 1 & 0,07 & 0,06 & 0,33 \\ \hline
Q2, R=20 & A & 7,5 & 4,0 & 40 & 0,61 & 0,39 & 1,10 \\ \hline
 & B & 1,0 & 1,0 & 20 & 0,08 & 0,06 & 0,33 \\ \hline
Q2R, R=10 & A & 6,0 & 4,0 & 20 & 0,45 & 0,29 & 1,13 \\ \hline
 & B & 1,0 & 1,0 & 10 & 0,12 & 0,07 & 0,43 \\ \hline
Q2R, R=10.000 & A & 10,5 & 10,5 & 19.485 & 101,12 & 5,04 & 12,50 \\ \hline
 & B & 6,0 & 6,0 & 9.742 & 4,28 & 0,28 & 2,18 \\ \hline
Q3, R=15.000 & A & 18,0 & 18,0 & 30.079 & 465,57 & 257,73 & 22,93 \\ \hline
 & B & 1,0 & 1,0 & 15.040 & 4,06 & 0,07 & 0,48 \\ \hline
\end{tabular}
\par\vspace{3pt}{\noindent\small\itshape Trung vị của 10 lệnh \code{EXPLAIN ANALYZE} mỗi lượt. Nguồn: results/rq1\_M2\_summary.csv.\par}
\end{table}

Trên TiDB, thiết kế không covering cần nhiều cop task và RPC hơn: Q2 dùng 7,5 cop task và 4 RPC, trong khi thiết kế covering chỉ dùng 1 cop task và 1 RPC. Với $R = 10$, phần ``còn lại'' chiếm 58\% thời gian RPC của thiết kế A, còn với $R = 10.000$ chỉ 11\% vì thời gian xử lý ở TiKV chiếm phần lớn. Đây là bằng chứng cho chi phí cố định ở $R$ nhỏ và chi phí mỗi dòng ở $R$ lớn, nhưng phần còn lại chỉ là xấp xỉ nên không được viết là ``chi phí mạng''. Thời gian RPC và xử lý cộng dồn qua các tác vụ song song, nên có thể lớn hơn độ trễ đo được (ví dụ Q2R, $R = 10.000$, thiết kế A: RPC 118,8 ms so với p50 44,3 ms). Riêng Q3 thiết kế A, TiDB ghi nhận 1.086 lần đọc block ngoài bộ nhớ đệm trên M2 và 20 lần trên M1, trong khi \code{Innodb\_buffer\_pool\_reads} của MySQL bằng 0; vì vậy một phần chênh lệch ở Q3 có thể đến từ việc đọc đĩa chứ không chỉ từ chi phí gọi KV.

\subsection{RQ1: đối chiếu giữa các máy}
Phép đo lặp lại ở cùng quy mô trên M1 (\path{results/rq1_M1_summary.csv}) cho cùng một dạng: ở $R = 100$, tỉ số là 1,41 (MySQL) và 2,15 (TiDB), so với 1,56 và 2,48 trên M2; ở $R = 10.000$ là 5,65 và 2,83, so với 5,91 và 3,31. Khác biệt duy nhất là Q3: trên M1 MySQL lợi hơn (12,09 so với 9,11), còn trên M2 hai hệ bằng nhau, nên không kết luận hệ nào lợi hơn ở Q3. Trên M3 (1 triệu dòng), $a_A / a_B$ là 1,04 ở MySQL và 1,53 ở TiDB, $b_A / b_B$ là 4,62 và 3,97, cùng chiều với M2 (2,07 của TiDB và 6,89 của MySQL); phép khớp của TiDB thiết kế A kém ($R^2 = 0{,}82$) vì điểm $R = 10.000$ dao động 42,8--91,1 ms. Vì các máy khác nhau, chỉ có thể nói dạng kết quả lặp lại trên cả ba máy; độ trễ tuyệt đối không được so.

\subsection{RQ2: ngưỡng độ chọn lọc}

\subsubsection{Thiết kế}
Câu hỏi là: khi truy vấn lấy một khoảng giá trị của \code{cod\_amount}, tỉ lệ dòng khớp bao nhiêu thì việc tra qua chỉ mục \code{idx\_cod} rồi quay về bảng bắt đầu chậm hơn quét cả bảng, và optimizer của mỗi hệ có đổi phương án đúng chỗ đó hay không. Nhóm không đặt trước một điểm giao nhau chung cho hai hệ.

Truy vấn có dạng \code{SELECT id, tracking\_code, status, cod\_amount, created\_at FROM shipments WHERE cod\_amount BETWEEN lo AND hi}, chạy ở ba chế độ: \emph{index} (ép dùng \code{idx\_cod} bằng \code{FORCE INDEX}), \emph{scan} (loại chỉ mục bằng \code{IGNORE INDEX}) và \emph{auto} (không gợi ý, để optimizer chọn). Có 10 mức độ chọn lọc: 0,001; 0,01; 0,1; 0,5; 1; 2; 5; 10; 25 và 50\% số dòng. Khoảng \code{[lo, hi]} của từng mức được tính từ seed 20261116 theo phân phối thực của cột, nên số dòng khớp trung bình $R$ nằm sát mục tiêu, từ 30 dòng (0,001\%) đến khoảng 1,5 triệu dòng (50\%); cột $R$ ở Bảng~\ref{tab:rq2-p50} là số dòng thực đo. Từ mức 2\% trở xuống, mỗi lượt chạy 20 khoảng khác nhau; từ mức 5\% trở lên chạy 3 khoảng vì mỗi truy vấn đã lâu. Mỗi ô chạy 5 lượt, tổng $2 \times 10 \times 3 \times 5 = 300$ lượt trên M3, hồ sơ XS, bảng 3 triệu dòng; trước mỗi lượt chương trình kiểm tra kế hoạch thực thi (đúng toán tử chỉ mục hoặc quét bảng) và số dòng trả về giữa ba chế độ. Các giá trị trong bảng là trung vị p50 của 5 lượt. Chỉ so hai hệ theo hình dạng đường cong và điểm giao nhau; không so độ trễ tuyệt đối với M1 hoặc M2.

\subsubsection{Kết quả}
\begin{table}[!htp]
\centering\bangnho
\caption{Độ trễ p50 (ms) theo độ chọn lọc và chế độ, 3 triệu dòng (M3)}\label{tab:rq2-p50}
\begin{tabular}{|C{1.3cm}|C{1.7cm}|C{1.45cm}|C{1.45cm}|C{1.45cm}|C{1.45cm}|C{1.45cm}|C{1.45cm}|}
\hline
\dongdau & & \multicolumn{3}{c|}{\textbf{MySQL}} & \multicolumn{3}{c|}{\textbf{TiDB C3}} \\ \hline
\dongdau \textbf{Mức} & \textbf{$R$} & \textbf{index} & \textbf{scan} & \textbf{auto} & \textbf{index} & \textbf{scan} & \textbf{auto} \\ \hline
0,001\% & 30 & 0,656 & 386 & 0,785 & 4,32 & 882 & 3,89 \\ \hline
0,01\% & 304 & 1,44 & 383 & 1,48 & 4,49 & 887 & 4,47 \\ \hline
0,1\% & 3.011 & 7,66 & 393 & 7,65 & 18,9 & 901 & 17,5 \\ \hline
0,5\% & 15.053 & 36,6 & 371 & 36,2 & 82,8 & 909 & 71,0 \\ \hline
1\% & 29.987 & 76,6 & 435 & 76,0 & 198 & 911 & 184 \\ \hline
2\% & 59.935 & 154 & 453 & 150 & 426 & 939 & 931 \\ \hline
5\% & 150.029 & 368 & 489 & 374 & 2.097 & 981 & 980 \\ \hline
10\% & 300.183 & 752 & 540 & 756 & 4.888 & 1.055 & 1.061 \\ \hline
25\% & 749.587 & 1.898 & 722 & 529 & 10.874 & 866 & 842 \\ \hline
50\% & 1.500.344 & 3.810 & 981 & 1.021 & 26.352 & 1.792 & 1.756 \\ \hline
\end{tabular}
\par\vspace{3pt}{\noindent\small\itshape Nguồn: results/rq2\_M3\_summary.csv. Trung vị của 5 lượt; ô auto của TiDB ở mức 2\% có kế hoạch lai (một phần truy vấn dùng chỉ mục, một phần quét bảng).\par}
\end{table}

Ở độ chọn lọc thấp, chỉ mục thắng rất rõ: ở 0,001\% MySQL cần 0,66 ms (index) so với 386 ms (scan), TiDB cần 4,3 ms so với 882 ms. Chi phí quét bảng gần như không đổi theo $R$ (MySQL khoảng 0,37--0,54 s, TiDB khoảng 0,88--1,06 s cho tới 10\%), vì đây là chi phí đọc cả bảng; chi phí tra qua chỉ mục tăng theo $R$. Hai đường cắt nhau ở độ chọn lọc trung bình, và tăng $R$ thêm nữa thì quét bảng thắng.

\begin{table}[!htp]
\centering\bangfont
\caption{Điểm giao nhau và điểm optimizer đổi phương án}\label{tab:rq2}
\begin{tabular}{|C{2.2cm}|C{3.0cm}|C{4.2cm}|C{5.2cm}|}
\hline
\dongdau \textbf{Hệ} & \textbf{Điểm giao nhau (\%)} & \textbf{Optimizer chuyển sang quét bảng} & \textbf{Hậu quả ở mức đo gần nhất} \\ \hline
MySQL & 6,9 (giữa 5 và 10) & Giữa 10 và 25\%; ở 10\% vẫn dùng chỉ mục & Ở 10\%: auto 756 ms, scan 540 ms (auto chậm hơn 1,4 lần) \\ \hline
TiDB C3 & 3,2 (giữa 2 và 5) & Từ 2\% (kế hoạch lai), hoàn toàn ở 5\% & Ở 2\%: auto 931 ms, index 426 ms (auto chậm hơn 2,2 lần) \\ \hline
\end{tabular}
\par\vspace{3pt}{\noindent\small\itshape Nguồn: results/rq2\_M3\_summary.csv. Điểm giao nhau nội suy tuyến tính trong thang log giữa hai mức đo liên tiếp; vì lưới thưa nên chỉ khoảng trong ngoặc là chắc chắn.\par}
\end{table}

\textbf{Điểm giao nhau.} MySQL giao ở khoảng 6,9\%: ở 5\% index nhanh hơn scan (368 so với 489 ms), ở 10\% chậm hơn (752 so với 540 ms). TiDB giao sớm hơn, khoảng 3,2\%: ở 2\% index nhanh hơn (426 so với 939 ms), ở 5\% chậm hơn nhiều (2.097 so với 981 ms). Với cả hai hệ, khoảng giữa hai mức đo không chồng lên khoảng nhỏ nhất--lớn nhất của các lượt, nên việc điểm giao nằm trong khoảng đó không phụ thuộc vào nhiễu; con số 6,9\% và 3,2\% chỉ là nội suy trong khoảng ấy.

\textbf{Optimizer.} Hai hệ lệch khỏi điểm giao nhau theo hai chiều ngược nhau. MySQL giữ chỉ mục quá lâu: ở 10\% optimizer vẫn chọn chỉ mục dù quét bảng nhanh hơn, và chỉ đổi sang quét bảng ở mức 25\% (đo được trong khoảng từ 10 đến 25\%). TiDB đổi quá sớm: ở 2\%, mức mà chỉ mục còn nhanh gấp 2,2 lần, kế hoạch \code{auto} đã lẫn \code{IndexLookUp} với \code{TableReader} ở cả 5 lượt, và p50 của \code{auto} (931 ms) gần bằng \code{scan} (939 ms) thay vì \code{index} (426 ms). Khoảng nhỏ nhất--lớn nhất của \code{auto} và \code{index} ở mức này không chồng nhau (706--941 so với 412--595 ms). Số liệu không cho biết nguyên nhân: nhóm không đo thống kê mà optimizer dùng để ước lượng số dòng, nên không kết luận optimizer ước lượng sai chỗ nào.

\textbf{Chi phí cố định và chi phí trên mỗi dòng.} Khớp mô hình $t = a + bR$ với 10 mức (trọng số $1/t$, \path{results/rq2_M3_fit.csv}) cho: MySQL index $a = 0{,}60$ ms, $b = 2{,}49$ $\mu$s/dòng; MySQL scan $a = 397$ ms, $b = 0{,}42$ $\mu$s/dòng, R$^2$ = 0,98; TiDB index $a = 3{,}05$ ms, $b = 7{,}17$ $\mu$s/dòng; TiDB scan $a = 890$ ms, $b = 0{,}30$ $\mu$s/dòng. Với MySQL, mô hình khớp tốt (R$^2$ = 0,9995 cho index): chi phí khoảng 2,5 $\mu$s mỗi dòng, hầu như không đổi trên toàn dải $R$ (từ 2,4 đến 2,6 $\mu$s). Với TiDB, mô hình tuyến tính khớp kém (R$^2$ = 0,56 cho index và 0,60 cho scan, phần dư lớn nhất 59\%). Lý do đọc được từ số liệu là chi phí trên mỗi dòng của TiDB không cố định: thời gian chia cho $R$ tăng từ khoảng 6,6 $\mu$s ở 1\% lên 14--18 $\mu$s từ 5\% trở lên. Vì vậy hệ số $b$ của TiDB (7,2 $\mu$s) chỉ là giá trị trung bình thô, và bảng này không được dùng để dự đoán $R$ nằm ngoài các mức đã đo. Nhóm chưa tách được nguyên nhân của sự tăng này; các bộ đếm chi tiết được lưu trong \path{results/rq2_M3.csv}.

\hinh{rq2_M3.png}{0.95}{Độ trễ p50 theo độ chọn lọc, thang log cả hai trục, ba chế độ; đường dọc là điểm giao nhau của từng hệ (M3, 3 triệu dòng)}{fig:rq2}

\textbf{Độ tin cậy.} Quét bảng của TiDB dao động nhiều giữa các lượt (656--1.139 ms ở cùng một mức, so với 359--389 ms của MySQL). Ở mức 25\%, hai chế độ của MySQL chạy cùng kế hoạch quét bảng nhưng cho 722 và 529 ms, nên chênh lệch dưới khoảng một phần ba được coi là nhiễu và kết luận chỉ dựa trên các khác biệt lớn hơn. Điểm giao nhau có thể dịch chuyển khi dữ liệu nhỏ hơn hoặc lớn hơn bộ nhớ đệm.

\subsection{Kết luận chương}
\textbf{RQ1: covering index và chi phí đọc.} Giả thuyết được xác nhận một phần. Ở $R$ nhỏ, covering lợi cho TiDB hơn (Q2R, $R = 10$: 2,01 so với 1,08 trên M2; 1,72 so với 1,07 trên M3), vì quay về bảng thêm vào TiDB khoảng 2,0 ms cố định so với 0,03 ms ở MySQL. Ở $R$ lớn, hai hệ không gần nhau mà MySQL lợi hơn trên M2 (5,91 so với 3,31 khi $R = 10.000$; trên M3 là 4,03 so với 3,80), có thể vì mỗi dòng quay về bảng ở MySQL là một lần tra B+Tree; điều này chưa được kiểm chứng. Ở Q3 không kết luận được hệ nào lợi hơn.

\textbf{RQ2: ngưỡng độ chọn lọc.} Điểm giao nhau giữa dùng chỉ mục và quét bảng là khoảng 6,9\% ở MySQL và 3,2\% ở TiDB (M3, 3 triệu dòng). Optimizer lệch theo hai chiều ngược nhau: MySQL giữ chỉ mục quá lâu (ở 10\%, chậm hơn quét bảng 1,4 lần), còn TiDB bỏ chỉ mục quá sớm (ở 2\%, chậm hơn dùng chỉ mục 2,2 lần).



\clearpage
\section{Chi phí ghi khi tăng số chỉ mục phụ}
\label{sec:ghi}
% Mục B.3. Phụ trách: TV2.
\phutrach{TV2 -- mục B.3}

\subsection{Câu hỏi và hai giả thuyết}
\ghichu{(a) LSM-tree làm TiDB giảm ít hơn; (b) cặp KV qua prewrite, commit và Raft làm TiDB giảm nhiều hơn. Cách tách: MySQL $\to$ C1 $\to$ C3; W1 và W2; covering và không covering.}

\subsection{Thiết kế thí nghiệm}
Thí nghiệm chạy trên máy M2 với 1 triệu dòng, 64 luồng, 60 giây mỗi lượt và 3 vòng. W1 thêm đơn mới; W2 cập nhật \code{status} và \code{updated\_at} theo \code{id}.

\begin{table}[!htp]
\centering\bangfont
\caption{Ma trận 84 lượt của RQ4}\label{tab:rq4-ma-tran}
\begin{tabular}{|L{4.2cm}|L{6.6cm}|C{3.4cm}|}
\hline
\dongdau \textbf{Cấu hình} & \textbf{Bộ chỉ mục} & \textbf{Số lượt} \\ \hline
MySQL & 0, 1, 3, 5, sender\_created, sender\_cov & $6 \times 2 \times 3 = 36$ \\ \hline
TiDB C3 & 0, 1, 3, 5, sender\_created, sender\_cov & $6 \times 2 \times 3 = 36$ \\ \hline
TiDB C1 & 0, 5 & $2 \times 2 \times 3 = 12$ \\ \hline
\end{tabular}
\end{table}
\ghichu{Nêu: bảng mới cho mỗi lượt, chờ đĩa lắng, kiểm tra bản sao, thứ tự xáo theo vòng, vì sao bỏ 16 luồng. W2 chỉ làm thay đổi chỉ mục có chứa \code{status}. Nguồn: \path{bench/rq4.py}, \path{scripts/rq4_official.ps1}, commit 317f0a6.}

\subsection{Kết quả thêm đơn mới (W1)}
\begin{table}[!htp]
\centering\bangfont
\caption{W1: TPS trung vị theo số chỉ mục và thay đổi so với 0 chỉ mục}\label{tab:rq4-w1}
\begin{tabular}{|C{2.3cm}|C{2.3cm}|C{2.3cm}|C{2.3cm}|C{2.3cm}|C{2.2cm}|}
\hline
\dongdau \textbf{Cấu hình} & \textbf{0} & \textbf{1} & \textbf{3} & \textbf{5} & \textbf{0 $\to$ 5} \\ \hline
MySQL & \chuaco & \chuaco & \chuaco & \chuaco & \chuaco \\ \hline
TiDB C3 & \chuaco & \chuaco & \chuaco & \chuaco & \chuaco \\ \hline
TiDB C1 & \chuaco & --- & --- & \chuaco & \chuaco \\ \hline
\end{tabular}
\par\vspace{3pt}{\noindent\small\itshape Nguồn: results/rq4\_official\_summary.csv. Trung vị của 3 lượt.\par}
\end{table}

\hinh{rq4_tps_vs_indexes.png}{0.95}{TPS theo số chỉ mục, W1 và W2, 64 luồng}{fig:rq4}

\subsection{Kết quả cập nhật trạng thái (W2)}
\begin{table}[!htp]
\centering\bangfont
\caption{W2: TPS trung vị theo số chỉ mục và thay đổi so với 0 chỉ mục}\label{tab:rq4-w2}
\begin{tabular}{|C{2.3cm}|C{2.3cm}|C{2.3cm}|C{2.3cm}|C{2.3cm}|C{2.2cm}|}
\hline
\dongdau \textbf{Cấu hình} & \textbf{0} & \textbf{1} & \textbf{3} & \textbf{5} & \textbf{0 $\to$ 5} \\ \hline
MySQL & \chuaco & \chuaco & \chuaco & \chuaco & \chuaco \\ \hline
TiDB C3 & \chuaco & \chuaco & \chuaco & \chuaco & \chuaco \\ \hline
TiDB C1 & \chuaco & --- & --- & \chuaco & \chuaco \\ \hline
\end{tabular}
\par\vspace{3pt}{\noindent\small\itshape Nguồn: results/rq4\_official\_summary.csv. Trung vị của 3 lượt.\par}
\end{table}

\subsection{Chi phí ghi của covering index}
\begin{table}[!htp]
\centering\bangfont
\caption{TPS với chỉ mục không covering và covering}\label{tab:rq4-covering}
\begin{tabular}{|C{2.4cm}|C{1.8cm}|C{3.4cm}|C{3.4cm}|C{2.8cm}|}
\hline
\dongdau \textbf{Cấu hình} & \textbf{Tải} & \textbf{sender\_created} & \textbf{sender\_cov} & \textbf{Thay đổi} \\ \hline
MySQL & W1 & \chuaco & \chuaco & \chuaco \\ \hline
MySQL & W2 & \chuaco & \chuaco & \chuaco \\ \hline
TiDB C3 & W1 & \chuaco & \chuaco & \chuaco \\ \hline
TiDB C3 & W2 & \chuaco & \chuaco & \chuaco \\ \hline
\end{tabular}
\end{table}

\subsection{Chi phí nhân bản: C1 và C3}
\ghichu{C1 và C3 ở bộ 0 và 5. Đây là chi phí nhân bản Raft trên một máy và một ổ đĩa, không phải chi phí mạng. Nêu leader của C1 lệch.}

\subsection{Độ trễ, byte ghi và tài nguyên}
\begin{table}[!htp]
\centering\bangnho
\caption{p95, p99, byte ghi mỗi giao dịch và CPU của sysbench}\label{tab:rq4-chi-tiet}
\begin{tabular}{|C{1.9cm}|C{1.2cm}|C{1.7cm}|C{1.9cm}|C{1.9cm}|C{2.5cm}|C{2.3cm}|}
\hline
\dongdau \textbf{Cấu hình} & \textbf{Tải} & \textbf{Bộ chỉ mục} & \textbf{p95 (ms)} & \textbf{p99 (ms)} & \textbf{Byte ghi / giao dịch} & \textbf{CPU sysbench (\%)} \\ \hline
MySQL & W1 & 0 / 5 & \chuaco & \chuaco & \chuaco & \chuaco \\ \hline
MySQL & W2 & 0 / 5 & \chuaco & \chuaco & \chuaco & \chuaco \\ \hline
TiDB C3 & W1 & 0 / 5 & \chuaco & \chuaco & \chuaco & \chuaco \\ \hline
TiDB C3 & W2 & 0 / 5 & \chuaco & \chuaco & \chuaco & \chuaco \\ \hline
TiDB C1 & W1 & 0 / 5 & \chuaco & \chuaco & \chuaco & \chuaco \\ \hline
TiDB C1 & W2 & 0 / 5 & \chuaco & \chuaco & \chuaco & \chuaco \\ \hline
\end{tabular}
\end{table}
\ghichu{TiKV bận nhất ở C3 thường dùng 0,7--0,8 trên 0,8 CPU, MySQL khoảng 3,9 trên 4 CPU. Dung lượng chỉ mục TiDB trong CSV là ước tính thống kê.}

\subsection{Các điểm bất thường}
\ghichu{Nêu đủ danh sách trong \path{docs/checklist_tv2.md}. Không bỏ lượt nào; bỏ thử 6 lượt bị nhiễu thì hướng kết luận không đổi.}

\subsection{Thảo luận}
\ghichu{Giả thuyết nào được số liệu ủng hộ, điểm nào chưa giải thích được. Phân biệt số đo và cách giải thích. Không quy chênh lệch cho mạng.}

\subsection{Kết luận chương}
\ghichu{Ba đến bốn câu, mỗi câu kèm số.}


\clearpage
\section{Tổng hợp, khuyến nghị và hạn chế}
\label{sec:tong-hop}
% Mục B.4, giả thuyết và khuyến nghị. Phụ trách: TV1.
\phutrach{TV1 -- mục B.4, bảng giả thuyết và khuyến nghị}

\subsection{Giả thuyết và kết quả}
\begin{table}[!htp]
\centering\bangnho
\caption{Đối chiếu giả thuyết với kết quả}\label{tab:gia-thuyet}
\begin{tabular}{|C{1.4cm}|L{4.0cm}|L{3.0cm}|L{3.8cm}|L{2.2cm}|}
\hline
\dongdau \textbf{Câu hỏi} & \textbf{Giả thuyết} & \textbf{Kết luận (xác nhận, bác bỏ, một phần)} & \textbf{Số liệu chính} & \textbf{Nguồn} \\ \hline
B.1 & \chuaco & \chuaco & \chuaco & \chuaco \\ \hline
RQ1 & \chuaco & \chuaco & \chuaco & \chuaco \\ \hline
RQ2 & \chuaco & \chuaco & \chuaco & \chuaco \\ \hline
RQ4 & \chuaco & \chuaco & \chuaco & \chuaco \\ \hline
\end{tabular}
\end{table}

\subsection{Đánh đổi của covering index}
\begin{table}[!htp]
\centering\bangfont
\caption{Lợi ích đọc, chi phí ghi và dung lượng của covering index}\label{tab:danh-doi}
\begin{tabular}{|L{5.6cm}|C{4.2cm}|C{4.2cm}|}
\hline
\dongdau \textbf{Chỉ số} & \textbf{MySQL} & \textbf{TiDB C3} \\ \hline
Lợi ích đọc Q2 (Bảng~\ref{tab:covering}) & \chuaco & \chuaco \\ \hline
Thay đổi TPS W1 (Bảng~\ref{tab:rq4-covering}) & \chuaco & \chuaco \\ \hline
Thay đổi TPS W2 (Bảng~\ref{tab:rq4-covering}) & \chuaco & \chuaco \\ \hline
Dung lượng thêm (Bảng~\ref{tab:dl-covering}) & \chuaco & \chuaco \\ \hline
\end{tabular}
\end{table}
\ghichu{Ba nguồn số đo trên ba máy khác nhau, nên chỉ ghép tỉ lệ, không ghép số tuyệt đối.}

\subsection{Khuyến nghị thiết kế chỉ mục}
\ghichu{Quy tắc nào giữ nguyên, quy tắc nào thay đổi khi chuyển từ MySQL sang TiDB. Mỗi khuyến nghị dẫn tới một bảng hoặc hình.}

% Mục B.4, hạn chế. Phụ trách: TV3.
\subsection{Hạn chế}
\phutrach{TV3 -- mục B.4, hạn chế}
Các kết luận của Chương~\ref{sec:doc} và \ref{sec:ghi} chỉ có giá trị trong phạm vi sau.
\begin{itemize}
  \item \textbf{Một máy, một ổ đĩa.} Mọi node của TiDB và MySQL chạy trong container trên cùng một máy và dùng chung một ổ đĩa. Giữa các container không có độ trễ mạng vật lý, nên chi phí gọi KV đo được chỉ gồm tuần tự hóa, gRPC và lập lịch; trên cụm thật, độ trễ mạng cộng thêm vào chi phí cố định $a$, và giá trị $a$ ở đây là ước lượng thấp. Độ trễ phía máy khách còn gồm đường từ máy chủ Windows vào container, giống nhau cho cả hai hệ.
  \item \textbf{Chịu lỗi.} Cấu hình C3 chịu mất một container TiKV, không phải mất một máy hay một ổ đĩa. Cấu hình C1 (3 TiKV, 1 bản sao) không chịu lỗi và chỉ dùng để tách chi phí nhân bản.
  \item \textbf{Tài nguyên không đối xứng.} CPU chia cố định theo container (hồ sơ XS: 4 CPU cho mỗi hệ). Tổng RAM của hai hệ khác nhau: cụm TiDB khoảng 5,5 GiB (5.632 MiB), MySQL 2 GiB, và bộ nhớ đệm đặt tương đương 750 MiB. Hai hệ cũng khác nhau về bộ thực thi, định dạng dòng và ngôn ngữ cài đặt, nên chênh lệch tuyệt đối giữa hai hệ không quy được cho riêng mạng hay riêng covering.
  \item \textbf{Ba máy khác nhau.} Mỗi máy có CPU, ổ đĩa và nền hệ điều hành riêng, nên không so độ trễ hay thông lượng tuyệt đối giữa các máy; chỉ so các đại lượng không có đơn vị (tỉ số covering, điểm giao nhau). Máy tính cá nhân không được cách ly hoàn toàn, nên có nhiễu ngoài tầm kiểm soát.
  \item \textbf{Quy mô.} Mọi máy dùng hồ sơ XS. RQ1 ở 1 triệu (M3) và 3 triệu dòng (M1, M2); RQ2 ở 3 triệu dòng; RQ4 chỉ ở 1 triệu dòng, 64 luồng, 3 vòng. \textbf{Không có số đo 10 triệu dòng}; mọi số liệu ở quy mô đó chỉ là số theo mô hình dung lượng. Ở mức chọn lọc trên 2\%, RQ2 chỉ dùng vài khoảng giá trị mỗi mức để vừa thời gian.
  \item \textbf{Bộ nhớ đệm kết quả của TiDB bị tắt} trong lúc đo RQ1 và RQ2 (MySQL 8 không có cơ chế tương đương); phép đo phản ánh đường truy cập chứ không phản ánh bộ nhớ đệm kết quả. Thời gian trong \code{EXPLAIN ANALYZE} cộng dồn qua các tác vụ song song, nên có thể lớn hơn thời gian chạy; phần còn lại \code{rest\_ms} chỉ là xấp xỉ.
  \item \textbf{Dữ liệu tổng hợp.} Dữ liệu vận đơn sinh bằng một bộ sinh cố định (seed 20261116), không phải dữ liệu sản xuất; phân bố \code{sender\_id} theo luật lũy thừa và \code{cod\_amount} phân bố đều.
\end{itemize}
Các điểm sau trên M3 có thể ảnh hưởng tới việc đọc kết quả; nhóm không xóa lượt nào.
\begin{itemize}
  \item \textbf{Số lượt.} RQ1 (1 triệu dòng) có 170 lượt, 85 mỗi hệ; RQ2 (3 triệu dòng) có 300 lượt. Trong cả 470 lượt, điều kiện lắng đĩa (\code{settled}) đều bằng 1; 17 lượt của TiDB trong RQ1 phải chờ 30,2 giây thay vì 20,1 giây.
  \item \textbf{Cơ sở dữ liệu của RQ2.} RQ2 trên M3 chạy trên cơ sở dữ liệu \code{logistics}, không phải \code{logistics\_3000000} như quy ước của nhóm. Bảng 3 triệu dòng ở hai nơi giống hệt nhau (tổng \code{cod\_amount} bằng 15001281300461), nên không ảnh hưởng tới số đo; bảng ở \code{logistics} được xóa sau đó để nạp 1 triệu dòng cho RQ1.
  \item \textbf{Lần chạy RQ1 đầu tiên dừng ở bước nạp dữ liệu} vì \code{logistics} đang chứa bảng 3 triệu dòng; chưa có phép đo nào được ghi, và lần chạy lại dùng đúng tham số của script chính thức.
  \item \textbf{Dao động giữa các lượt:} Q3 và Q2R $R = 10.000$ của TiDB trong RQ1 (mục B.2), quét bảng của TiDB trong RQ2.
\end{itemize}


% ===== 9. PHẦN KẾT =====
% Kết luận. Phụ trách: TV1.
\clearpage
\section*{Kết luận}
\addcontentsline{toc}{section}{Kết luận}
\phutrach{TV1}

\subsection*{Kết quả đạt được}
\addcontentsline{toc}{subsection}{Kết quả đạt được}

Đề tài đã hoàn thành các phép so sánh chính về dung lượng, chi phí đọc và chi phí ghi trên dữ liệu vận đơn và hai hệ TiDB v8.5.5/MySQL 8.4.11 trong phạm vi môi trường lab đã ghi. Về cấu trúc, cụm TiDB gồm 1 PD, 1 TiDB Server và 3 TiKV; MySQL chạy một node InnoDB. Mọi kết quả tuyệt đối được giới hạn trong máy, hồ sơ tài nguyên và số lượt đo tương ứng.

\begin{enumerate}
  \item \textbf{Mô hình và dung lượng.} Mô hình dự đoán dung lượng trang lá của MySQL lệch không quá 1,2\%, còn mô hình TiDB khớp dữ liệu và các index thường với sai lệch khoảng 0--4\%. Riêng \code{idx\_code} UNIQUE của TiDB lớn hơn mô hình 24--32 B/mục và nguyên nhân giá trị dư chưa được xác minh đầy đủ. Mô hình cho bộ 5 ở 10 triệu dòng là khoảng 2.649 MiB trên MySQL và 5.598 MiB trên một bản sao TiDB; đây chỉ là kết quả ngoại suy, không phải số đo. Chi tiết ở Chương~\ref{sec:dung-luong}.

  \item \textbf{RQ1 -- covering và quay về bảng.} Trên M2 ở 3 triệu dòng, covering làm truy vấn Q2R với $R=10$ nhanh gấp 2,01 lần trên TiDB C3 và 1,08 lần trên MySQL. Khi $R=10.000$, tỷ lệ tương ứng là 3,31 và 5,91 lần. Kết quả chứng minh covering có lợi lớn ở dải $R$ nhỏ, nhưng lợi ích tương đối không tỷ lệ thuận với số dòng. Tại $R$ lớn, MySQL có lợi hơn trong phép đo M2.

  \item \textbf{RQ2 -- độ chọn lọc và optimizer.} Điểm giao \emph{performance crossover} là khoảng 6,9\% ở MySQL và 3,2\% ở TiDB C3 trên M3, ba triệu dòng. Optimizer của MySQL giữ index ở mức 10\% dù quét bảng nhanh hơn khoảng 1,4 lần; TiDB có kế hoạch hỗn hợp giữa \code{IndexLookUp} và \code{TableReader} tại mức 2\% dù dùng index còn nhanh gấp 2,2 lần. Hai điểm crossover và optimizer crossover không trùng nhau ở cả hai hệ.

  \item \textbf{RQ3 -- kiểu khóa chính.} Khóa \code{tracking\_code} làm \code{idx\_sender\_created} tăng khoảng 27,4\% trên MySQL và 39,5\% trên TiDB so với \code{id} ở 1 triệu dòng. TiDB clustered \code{id} tra bằng một gọi TiKV; \code{id} nonclustered và khóa chuỗi cần thêm một pha quay về dòng. Khóa chuỗi clustered có thể tra mã nhanh nhất nhưng đánh đổi dung lượng cho mọi chỉ mục phụ.

  \item \textbf{RQ4 -- workload ghi.} LSM-Tree không làm mức giảm TPS của TiDB nhỏ hơn MySQL trong thí nghiệm. Từ 0 tới 5 chỉ mục, TiDB C3 giảm TPS 82,7\% ở W1 và 64,8\% ở W2; MySQL giảm 18,9\% và 20,2\%. C1 vẫn giảm 79,4\% ở W1 và 51,6\% ở W2, nên phần lớn chi phí không đến riêng từ nhân bản Raft ba bản sao. Kết quả bác bỏ giả thuyết rằng TiDB giảm ít hơn nhờ ghi tuần tự LSM-Tree; giả thuyết chi phí giao dịch phân tán giữ vai trò chính được số đo hiện hành hợp lý.
\end{enumerate}

\subsection*{Hạn chế}
\addcontentsline{toc}{subsection}{Hạn chế}
Các con số hiện tại không chứng minh một hệ luôn tốt hơn hệ kia. Mọi TiDB/TiKV và MySQL đều chạy trong container trên một máy và một ổ đĩa, nên không đo độ trễ mạng vật lý đa máy. RQ4 chỉ đo 1 triệu dòng, 64 luồng và 3 vòng; RQ1 đo 1 và 3 triệu dòng; RQ2 đo 3 triệu dòng; RQ3 chỉ đo 1 triệu dòng. Dữ liệu là dữ liệu tổng hợp, không phải log vận đơn sản xuất. Các giới hạn chi tiết gồm thiếu số đo chính thức ở 10 triệu dòng, cấu hình cache và bộ máy khác nhau, khác biệt collation đã kiểm tra và nhiễu cửa sổ 60 giây đã ghi trong Chương~\ref{sec:tong-hop}.

\subsection*{Hướng phát triển}
\addcontentsline{toc}{subsection}{Hướng phát triển}

\begin{enumerate}
  \item Chạy lại RQ1--RQ4 trên ít nhất ba node vật lý riêng, đo độ trễ mạng và I/O đĩa độc lập; xác nhận cấu hình CPU/RAM/cache bằng cùng hồ sơ.
  \item Mở rộng quy mô lên 10 triệu và 30 triệu dòng, bổ sung số đo dung lượng sau nén trên đĩa, TiKV MVCC và compaction thay vì chỉ mô hình logic.
  \item Chạy RQ2 với lưới độ chọn lọc dày hơn quanh 2--10\%, lưu thống kê optimizer để phân biệt chi phí ước lượng sai và chi phí cơ chế thực thi.
  \item Thêm các workload hỗn hợp đọc--ghi và mức luồng khác nhau; dùng \code{run\_id}, raw artifact và thống kê multi-run để kiểm định độ tin cậy.
  \item Với TiDB, đánh giá TiFlash cho truy vấn phân tích nhưng tách rõ chi phí giao dịch OLTP và lưu trữ HTAP; với MySQL, so thêm InnoDB Cluster hoặc sharding chỉ khi dữ liệu/tính sẵn sàng được thiết kế ngang hàng.
\end{enumerate}

\subsection*{Kiến thức và kỹ năng đạt được}
\addcontentsline{toc}{subsection}{Kiến thức và kỹ năng đạt được}

\textbf{TV1} hiểu cách liên kết mô hình lưu trữ, truy cập khóa chính, chỉ mục phụ và kết quả thực nghiệm; rèn kỹ năng tạo mô hình dung lượng, kiểm tra byte và đối chiếu B+Tree/LSM-Tree.

\textbf{TV2} nắm quy trình dựng cụm Docker, cấu hình tài nguyên, nạp dữ liệu, chạy benchmark ghi, trích cấu hình thực tế và xử lý các lượt đo bất thường; kỹ năng này giúp tách điều kiện benchmark khỏi suy diễn về hiệu năng.

\textbf{TV3} hiểu cách phân tích \code{EXPLAIN ANALYZE}, mapping chỉ mục trên TiDB/MySQL, selectivity crossover và thống kê latency; rèn khả năng đọc số liệu, giữ raw artifact và không loại lượt đo chỉ vì kết quả trái giả thuyết.

% ==============================================================================
% TÀI LIỆU THAM KHẢO (THEO CHUẨN TRÍCH DẪN KHOA HỌC)
% ==============================================================================

\clearpage
\phantomsection
\addcontentsline{toc}{section}{\refname}
\begin{thebibliography}{99}

\bibitem{pingcap_arch}
PingCAP, \textit{TiDB Architecture Overview}, PingCAP Documentation, 2024. [Online]. Available: \url{https://docs.pingcap.com/tidb/stable/tidb-architecture}.

\bibitem{pingcap_mapping}
PingCAP, \textit{Mapping Table Data to Key-Value Pairs in TiDB}, PingCAP Engineering Guides, 2024. [Online]. Available: \url{https://docs.pingcap.com/tidb/stable/tidb-computing}.

\bibitem{tikv_storage}
TiKV Authors, \textit{TiKV: A Distributed Transactional Key-Value Database}, CNCF Graduated Project Documentation, 2024. [Online]. Available: \url{https://tikv.org/docs/}.

\bibitem{mysql_innodb}
Oracle Corporation, \textit{MySQL 8.4 Reference Manual: InnoDB Physical Structure and Multi-Versioning}, Oracle Database Documentation, 2024. [Online]. Available: \url{https://dev.mysql.com/doc/refman/8.4/en/innodb-storage-engine.html}.

\bibitem{spanner}
J. C. Corbett et al., ``Spanner: Google's Globally-Distributed Database,'' in \textit{Proceedings of the 10th USENIX Conference on Operating Systems Design and Implementation (OSDI '12)}, Hollywood, CA, USA, 2012, pp. 251--264.

\bibitem{f1}
J. Shute et al., ``F1: A Distributed SQL Database That Scales,'' in \textit{Proceedings of the VLDB Endowment}, vol. 6, no. 11, pp. 1068--1079, 2013.

\bibitem{percolator}
D. Peng and F. Dabek, ``Large-scale Incremental Processing Using Distributed Transactions and Notifications,'' in \textit{Proceedings of the 9th USENIX Conference on Operating Systems Design and Implementation (OSDI '10)}, Vancouver, BC, Canada, 2010, pp. 251--264.

\bibitem{raft}
D. Ongaro and J. Ousterhout, ``In Search of an Understandable Consensus Algorithm,'' in \textit{2014 USENIX Annual Technical Conference (USENIX ATC 14)}, Philadelphia, PA, USA, 2014, pp. 305--319.

\bibitem{lsmtree}
P. O'Neil, E. O'Neil, and G. Weikum, ``The Log-Structured Merge-Tree (LSM-Tree),'' \textit{Acta Informatica}, vol. 33, no. 4, pp. 351--385, 1996.

\bibitem{rocksdb}
Meta Platforms / The RocksDB Team, \textit{RocksDB: A Persistent Key-Value Store for Fast Storage Environments}, Facebook Open Source, 2024. [Online]. Available: \url{https://rocksdb.org/}.

% TV1: chưa xác minh được nguồn này, kiểm tra hoặc bỏ trước khi nộp.
\bibitem{pipedrive}
Pipedrive Engineering, \textit{Evaluating TiDB as a Distributed Replacement for Large-Scale MySQL Workloads: An Empirical Study on 1.47 Billion Rows Table}, Technical Report, 2023.

\bibitem{sysbench}
A. Kopytov, \textit{Sysbench: Scriptable Multi-Threaded Benchmark Tool}, GitHub Repository, 2024. [Online]. Available: \url{https://github.com/akopytov/sysbench}.

\bibitem{db_concepts}
A. Silberschatz, H. F. Korth, and S. Sudarshan, \textit{Database System Concepts}, 7th ed. New York, NY, USA: McGraw-Hill Education, 2019.

\bibitem{tidb_sysvars}
PingCAP, \textit{System Variables} (mục \code{tidb\_enable\_1pc}), TiDB Documentation. [Online]. Available: \url{https://docs.pingcap.com/tidb/stable/system-variables/}. Truy cập ngày 24/11/2026.

\bibitem{tidb_hotspot}
PingCAP, \textit{Troubleshoot Hotspot Issues}, TiDB Documentation. [Online]. Available: \url{https://docs.pingcap.com/tidb/stable/troubleshoot-hot-spot-issues/}. Truy cập ngày 24/11/2026.

\bibitem{mysql_trx}
Oracle Corporation, \textit{MySQL 8.4 Reference Manual: Optimizing InnoDB Transaction Management}. [Online]. Available: \url{https://dev.mysql.com/doc/refman/8.4/en/optimizing-innodb-transaction-management.html}. Truy cập ngày 24/11/2026.

\bibitem{mysql_params}
Oracle Corporation, \textit{MySQL 8.4 Reference Manual: InnoDB Startup Options and System Variables}. [Online]. Available: \url{https://dev.mysql.com/doc/refman/8.4/en/innodb-parameters.html}. Truy cập ngày 24/11/2026.

\end{thebibliography}

% Phụ lục. Phụ trách: TV2 ghép, mỗi người thêm dòng của mình.
\clearpage
\section*{Phụ lục}
\addcontentsline{toc}{section}{Phụ lục}
\phutrach{TV2}

\subsection*{A. Lệnh chạy lại thí nghiệm}
\addcontentsline{toc}{subsection}{A. Lệnh chạy lại thí nghiệm}

Các lệnh chạy trong thư mục gốc của dự án, trên PowerShell. Nếu một lượt đo bị ngắt, chạy lại đúng lệnh đó để đo tiếp phần còn thiếu.

\begin{lstlisting}[caption={Lệnh PowerShell chạy lại các thí nghiệm}]
.venv\Scripts\python -m bench.cli check

# RQ1: 1 trieu dong (M3) va 3 trieu dong (M1)
.\scripts\rq1_official.ps1
.\scripts\rq1_official.ps1 -Rows 3000000
.venv\Scripts\python -m experiments.rq1_summary results/rq1_M3.csv

# RQ4: 84 luot (M2)
.\scripts\rq4_official.ps1
.venv\Scripts\python -m experiments.rq4_summary results/rq4_official.csv

# Demo chuc nang co ban
.\scripts\use_engine.ps1 -Engine mysql -Switch
.venv\Scripts\python scripts\demo_co_ban.py mysql
.\scripts\use_engine.ps1 -Engine tidb -Switch
.venv\Scripts\python scripts\demo_co_ban.py tidb
\end{lstlisting}
\ghichu{TV1 thêm lệnh B.1, TV3 thêm lệnh RQ2.}

\subsection*{B. Danh mục file kết quả}
\addcontentsline{toc}{subsection}{B. Danh mục file kết quả}
\begin{center}
\bangfont
\begin{tabular}{|C{1.8cm}|C{1.6cm}|L{7.4cm}|L{3.2cm}|}
\hline
\dongdau \textbf{Phần} & \textbf{Máy} & \textbf{File} & \textbf{Commit lúc đo} \\ \hline
B.1 & M1 & results/b1\_storage\_M1.csv & \chuaco \\ \hline
RQ1 & M1 & results/rq1\_M1.csv & \chuaco \\ \hline
RQ1 & M3 & results/rq1\_M3.csv & \chuaco \\ \hline
RQ2 & M3 & results/rq2\_M3.csv & \chuaco \\ \hline
RQ4 & M2 & results/rq4\_official.csv, results/rq4\_official\_summary.csv & 317f0a6 \\ \hline
\end{tabular}
\end{center}
\ghichu{TV1 và TV3 điền commit in ra ở dòng đầu khi chạy lệnh kiểm tra.}

\subsection*{C. Tài liệu bổ sung}
\addcontentsline{toc}{subsection}{C. Tài liệu bổ sung}
\begin{center}
\bangfont
\begin{tabular}{|L{5.4cm}|L{9.2cm}|}
\hline
\dongdau \textbf{Nội dung} & \textbf{Vị trí} \\ \hline
Mã nguồn, cấu hình, kết quả & \url{https://github.com/taesikwoo268/nhom7-tidb-indexing} \\ \hline
Hướng dẫn cài đặt và chạy & \path{README.md} \\ \hline
Cấu hình container & \path{docker/} \\ \hline
Chương trình đo dùng chung & \path{bench/} \\ \hline
Thí nghiệm RQ1, tổng hợp RQ4 & \path{experiments/} \\ \hline
Kịch bản tải W1, W2 & \path{sysbench/shipments_write.lua} \\ \hline
Demo chức năng cơ bản & \path{scripts/demo_co_ban.py} \\ \hline
Thiết kế kỹ thuật & \path{docs/2026-11-16-tidb-indexing-van-don.vi.md} \\ \hline
\end{tabular}
\end{center}
\ghichu{Trước khi nộp, kiểm tra người chấm mở được repo. Đầu ra gốc không nằm trong Git (\path{results/raw}) thì nộp kèm hoặc ghi nơi lưu.}


\end{document}
